\documentclass[11pt]{article}
\usepackage[T1]{fontenc}
\IfFileExists{lmodern.sty}{\usepackage{lmodern}}{}
\usepackage[margin=1in]{geometry}
\usepackage{amsmath,amssymb,amsthm,mathtools}
\usepackage{microtype}
\usepackage{enumitem}
\usepackage[colorlinks=true,linkcolor=blue,citecolor=blue,urlcolor=blue]{hyperref}
\hypersetup{pdftitle={Finitely presented infinite p-groups for every prime},pdfauthor={Seth Douglas and Nidhal Mghirbi}}

\numberwithin{equation}{section}
\newtheorem{theorem}{Theorem}[section]
\newtheorem{lemma}[theorem]{Lemma}
\newtheorem{proposition}[theorem]{Proposition}
\newtheorem{corollary}[theorem]{Corollary}
\newtheorem*{theoremA}{Theorem A}
\newtheorem*{theoremB}{Theorem B}
\newtheorem*{theoremC}{Theorem C}
\theoremstyle{definition}
\newtheorem{definition}[theorem]{Definition}
\theoremstyle{remark}
\newtheorem{remark}[theorem]{Remark}

\newcommand{\F}{\mathbb F}
\newcommand{\Q}{\mathbb Q}
\newcommand{\Z}{\mathbb Z}
\DeclareMathOperator{\St}{St}
\DeclareMathOperator{\E}{E}
\DeclareMathOperator{\GL}{GL}
\DeclareMathOperator{\SL}{SL}

\title{Finitely presented infinite $p$-groups for every prime}
\author{Seth Douglas \and Nidhal Mghirbi}
\date{October 2026}

\begin{document}
\maketitle

\begin{abstract}
For every prime $p$ we construct an infinite, residually finite group with an ordinary finite presentation in which every
element has finite $p$-power order. For $p=2$ this is a recent construction of OpenAI, which also answered the finitely
presented Burnside question. We extend it to every prime, using a single finite presentation written over the
integers. Over every field of positive characteristic this presentation defines a finitely presented,
infinite-dimensional, nil, Jacobson radical algebra that is not nilpotent and whose matrix algebras are all nil.
Reducing it modulo $p$ gives a graded $\F_p$-algebra $R_p$, and the groups are finite-index subgroups of the Steinberg
groups $\St_{12}(R_p)$. The
extension is not a relabelling of $2$ by $p$. The characteristic-two cancellation bound must be replaced by one that depends on
$p$. The unordered-frame homological stability argument fails over odd fields and is replaced
by an ordered-frame argument with signed transpositions. The phases that detect central elements must be taken modulo
integers coprime to $p$. Finally, we show that characteristic zero is genuinely different: over a field of characteristic zero, an algebra
presented by finitely many word equalities and word zeros is nil only if it is nilpotent. A characteristic-zero analogue
therefore needs relations with signs or coefficients.
\end{abstract}

\noindent\emph{The mathematics in this paper was developed with AI assistance; Section~\ref{sec:methods} describes how,
and the authors take responsibility for every claim.}

\setcounter{tocdepth}{1}
{\small\tableofcontents}

%======================================================================
\section{Introduction}\label{sec:intro}

A group is \emph{periodic} if each of its elements has finite order, and a \emph{$p$-group} (in the infinite sense used
here) if each element has finite $p$-power order. Burnside asked in 1902 whether periodic groups with finitely many
generators must be finite~\cite{Burnside1902}. Golod's nil algebras gave infinite finitely generated residually finite
$p$-groups for every prime~\cite{Golod1964,GolodShafarevich1964}, and Novikov and Adian gave infinite finitely generated
groups of bounded exponent~\cite{NovikovAdian1968}. A stronger requirement remained open: an \emph{ordinary finite
presentation}, meaning finitely many generators and finitely many relator words. Whether an infinite periodic group can
have one is recorded as an open question by Ol'shanskii and Sapir~\cite[p.~45]{OlshanskiiSapir2003}, and again by
Haettel and Osajda~\cite[Section~2.2(IX)]{HaettelOsajda2025}.

OpenAI recently answered this question~\cite{OAI247periodic}. They constructed a finitely presented graded algebra
$R=\F_2\oplus I$ over $\F_2$ in which every matrix with entries in $I$ is nilpotent, and showed that the Steinberg group
$\St_{12}(R)$ is infinite, finitely presented and periodic. A companion paper~\cite{OAI247rf} proved that $\St_{12}(R)$
is residually finite. The kernel of its reduction to $\St_{12}(\F_2)$ is then an infinite, finitely presented,
residually finite $2$-group. Both constructions are written in characteristic two. The finitely presented algebra
of~\cite{OAI247periodic} also answers, over $\F_2$, Ufnarovskij's question whether a finitely presented nil algebra must
be nilpotent~\cite[Question~8.1]{Smoktunowicz2008}, together with the finite-presentation versions of Amitsur's radical
question and Kurosh's algebraic question.

We extend these results to every prime, and we explain why characteristic zero lies outside the method.

\begin{theoremA}[Nil algebras in every positive characteristic]
There is a finite alphabet $\Sigma$ and a finite set of relations, each of the form $u=v$ with $u,v$ nonempty words of the
same length or $w=0$ with $w$ a nonempty word, such that the following holds for every field $K$ of positive
characteristic. Let $I_K$ be the nonunital $K$-algebra with this presentation. Then $I_K$ is finitely presented,
infinite-dimensional and not nilpotent, and for every $l\geq1$ every matrix in $M_l(I_K)$ is nilpotent. Consequently
$I_K$ is Jacobson radical, and its unitization $K\oplus I_K$ is a finitely presented, algebraic, infinite-dimensional
unital $K$-algebra.
\end{theoremA}

The presentation is fixed once, independently of $K$. For $K=\F_2$ the statement is~\cite{OAI247periodic}.

\begin{theoremB}[Finitely presented $p$-groups for every prime]
Let $p$ be a prime, and let $R_p=\F_p\oplus I_{\F_p}$ be the unitization of the algebra of Theorem~A over $\F_p$. Then:
\begin{enumerate}[label=\textup{(\roman*)}]
\item $\Gamma_p=\St_{12}(R_p)$ is infinite, ordinarily finitely presented, periodic and residually finite; and
\item $G_p=\ker\bigl(\Gamma_p\to\St_{12}(\F_p)\bigr)$ is an infinite, ordinarily finitely presented, residually finite
group in which every element has finite $p$-power order.
\end{enumerate}
\end{theoremB}

For $p=2$ this is~\cite{OAI247periodic,OAI247rf}. One fixed presentation thus gives the groups for every prime at once.
Two consequences follow from known results.
\begin{itemize}[leftmargin=*]
\item By Ershov and Jaikin-Zapirain~\cite[Theorem~6.2]{ErshovJaikinZapirain2010}, $\St_n(S)$ has Kazhdan's property~(T)
  for every finitely generated unital ring $S$ and $n\geq3$. So $\Gamma_p$ has property~(T), and so does its
  finite-index subgroup $G_p$. This is the odd-prime analogue of~\cite[Corollary~1.2]{OAI247periodic}.
\item The element orders of $G_p$ are unbounded. By Zel'manov's solution of the restricted Burnside
  problem~\cite{ZelmanovOdd,ZelmanovTwo}, a finitely generated residually finite group of exponent dividing $p^a$ is
  finite.
\end{itemize}

\begin{theoremC}[Characteristic zero needs signs]
Let $K$ be a field of characteristic zero, and let $A$ be a nonunital $K$-algebra generated by finitely many letters
subject only to relations $u=v$ and $w=0$, where $u,v,w$ are nonempty words. If $A$ is nil, then $A$ is nilpotent. In
particular, over every field of characteristic zero the presentation of Theorem~A defines an infinite-dimensional
algebra that is not nil.
\end{theoremC}

Theorem~C is elementary, and we do not present it as new (Section~\ref{sec:charzero}). We include it because it locates
the positive-characteristic mechanism precisely. The sum $x$ of the generators has $x^n$ equal to the sum of all words of length $n$. Its coefficient
on each surviving word class is a positive integer, which cannot vanish in characteristic zero. So $x^N=0$ forces every
word of length $N$ to vanish. Theorem~C shows exactly where the positive-characteristic mechanism lives: multiplicities
that are positive integers can vanish modulo $p$, never over $\Q$. A characteristic-zero finitely presented nil algebra
that is not nilpotent, if one exists, needs relations with signs or coefficients. Whether one exists remains open.

\paragraph{Why odd primes are not a relabelling.} Three steps of the characteristic-two argument do not survive the
change of field as written, and a fourth needs new parameters.
\begin{enumerate}[leftmargin=*]
\item \emph{Cancellation.} In characteristic two the symmetric sum $\sum_{\sigma\in S_m}X_{\sigma(1)}\cdots
  X_{\sigma(m)}$ of $d\times d$ matrices vanishes once $m>d^2$~\cite{OAI247periodic}. In characteristic $p$ this bound is false for small
  $d$, and the bound must depend on $p$. By a theorem of Zalesskii the sum vanishes once $m\geq dp$~\cite{Zalesskii1985}.
  We use the cruder bound $m>(p-1)d^2$, which has a two-line proof: some basis matrix then repeats at least $p$ times
  and contributes a factorial divisible by $p$. No signs are needed. The hierarchy absorbs the factor $p-1$ at a finite level (Section~\ref{sec:algebra}).
\item \emph{Homological stability.} The rational, unordered-frame stability argument of~\cite{OAI247periodic} relies on
  frame transpositions being elementary. Over an odd field this fails: we give an explicit stabilizer element whose
  projection is not elementary (Proposition~\ref{prop:unordered-fails}). We replace it with an integral ordered-frame
  argument built on the universal central extension of pointwise frame stabilizers, following~\cite{OAI247rf}. The key
  face-map comparison uses a \emph{signed} transposition that fixes the complementary block exactly
  (Section~\ref{sec:stability}).
\item \emph{Signs throughout the group calculations.} Row interchanges become signed transpositions. Diagonal units and
  permutations of used coordinates are corrected by a compensating sign on an unused coordinate. Torsion of reverse
  commutators uses the identity $H^p=[x^p,y]$ rather than an involution (Sections~\ref{sec:fp}--\ref{sec:torsion}).
\item \emph{Phases and permutations.} Detecting central elements in finite quotients uses long-word phases modulo
  integers coprime to the characteristic. In characteristic two the odd encoding lengths $L_j=2\cdot2^{h_j}+1$ serve
  directly. For odd $p$ we reduce the phases modulo the part of $L_0\cdots L_d$ prime to $p$. By the lifting-the-exponent
  lemma this is unbounded as $d\to\infty$, so one presentation serves every prime. The characteristic-two lemma that every
  coordinate permutation is induced on roots by conjugation~\cite[Lemma~3.3]{OAI247rf} is false for odd $p$. We replace
  it by signed transpositions with a compensating sign on an unused coordinate (Sections~\ref{sec:torsion}
  and~\ref{sec:detection}).
\end{enumerate}

\paragraph{What is due to whom.} The construction and the principal methods are OpenAI's~\cite{OAI247periodic,OAI247rf}:
\begin{itemize}[leftmargin=*]
\item the finite rule hierarchy and its fixed-point program;
\item matrix-weighted cancellation through cuts, tokens and weighted automata;
\item the coefficient extraction that proves matrix nilness;
\item the finite presentation of Steinberg groups (after Krsti\'c and McCool~\cite{KrsticMcCool});
\item root strips, centrality and stability;
\item the cyclic and truncated shift substitutions;
\item the long-word phases and the detection argument.
\end{itemize}
Their hierarchy is characteristic-free, and we use it unchanged. What is new here:
\begin{itemize}[leftmargin=*]
\item the extension to every prime, and the observation that one integral presentation serves every positive
  characteristic;
\item the failure of the unordered-frame port and its ordered-frame replacement with an exact signed face map;
\item the sign-corrected group calculations;
\item the reduction of phases to moduli prime to $p$, giving one presentation for every prime;
\item Theorem~C, which is elementary and included for completeness.
\end{itemize}

\paragraph{Organization.}
\begin{itemize}[leftmargin=*]
\item Section~\ref{sec:prelim}: word algebras and the inputs we import.
\item Section~\ref{sec:algebra}: Theorem~A.
\item Section~\ref{sec:charzero}: Theorem~C.
\item Sections~\ref{sec:fp}--\ref{sec:detection}: Theorem~B (finite presentation over $\F_p$, centrality and stability,
  relative $p$-primary torsion, then phases and detection).
\item Section~\ref{sec:discussion}: questions.
\item Section~\ref{sec:methods}: how the work was done and verified.
\end{itemize}

%======================================================================
\section{Preliminaries}\label{sec:prelim}

Throughout, $\Sigma$ denotes a finite alphabet, words are finite strings over $\Sigma$ (including the empty word), and
$|w|$ is the length of a word $w$. For a ring $R$ we write $[w]$ for the image of a word in an algebra presented on
$\Sigma$.

\subsection{Word presentations over the integers}

An \emph{elementary replacement} is a relation $u=v$ between nonempty words of the same length; it may be used in either
direction inside a word. A \emph{zero declaration} is a relation $w=0$ with $w$ nonempty. For a commutative ring
$\Lambda$ and a finite set $\mathcal S$ of elementary replacements and zero declarations, let $A(\Lambda)$ be the unital
associative $\Lambda$-algebra generated by $\Sigma$ subject to $\mathcal S$. Thus
$A(\Lambda)=\Lambda\langle\Sigma\rangle/\mathcal J_\Lambda$, where $\mathcal J_\Lambda$ is the two-sided ideal generated
by the binomials $u-v$ and the words $w$ listed in $\mathcal S$.

The \emph{replacement graph} has the words as vertices and an edge between $xuy$ and $xvy$ for every listed replacement
$u=v$ and all words $x,y$. A connected component of this graph is a \emph{replacement class}. A class is \emph{zero} if
one of its words contains a zero-declared word as a substring, and \emph{nonzero} otherwise. The following lemma is
\cite[Lemma~2.1]{OAI247periodic}, stated there over a field; we need it over $\Z$, so that a single presentation can be
used in every characteristic.

\begin{lemma}\label{lem:word-basis}
The images of the nonzero replacement classes form a basis of $A(\Z)$ as a free $\Z$-module, and the zero classes map to
$0$. For every commutative ring $\Lambda$, the natural map $\Lambda\otimes_{\Z}A(\Z)\to A(\Lambda)$ is an isomorphism, so
the images of the nonzero classes form a $\Lambda$-basis of $A(\Lambda)$. In particular, whether a word is zero does not
depend on $\Lambda$, and $A(\Lambda)$ is graded by word length with $A(\Lambda)_0=\Lambda$.
\end{lemma}

\begin{proof}
Let $V$ be the free $\Z$-module on all words, and let $V_1\subseteq V$ be the span of the differences $x-y$ for edges
$\{x,y\}$ of the replacement graph. A finite combination $\sum_w c_ww$ lies in $V_1$ if and only if, for every
replacement class $C$, the coefficient sum $\sum_{w\in C}c_w$ vanishes. Indeed, edge differences have this property. Conversely, within one class each word is
joined to a fixed base word by a path, and the difference of its endpoints telescopes into edge differences. Hence $V/V_1$
is free on the classes. The two-sided ideal of $\Z\langle\Sigma\rangle$ generated by the listed binomials is $V_1$, since
$V_1$ is spanned by the products $x(u-v)y$.

Multiplication of words on either side maps replacement classes to replacement classes, and maps zero classes to zero
classes. Let $V_0\subseteq V/V_1$ be the span of the zero classes. It is a two-sided ideal, it contains the images of the
zero-declared words, and it is contained in the ideal they generate, because every word of a zero class is equivalent,
modulo $V_1$, to a word $xwy$ with $w$ zero-declared. Hence $A(\Z)=(V/V_1)/V_0$ is free on the nonzero classes.

For an arbitrary $\Lambda$, right exactness of $\Lambda\otimes_{\Z}-$ applied to the presentation of $A(\Z)$ gives
$A(\Lambda)\cong\Lambda\otimes_{\Z}A(\Z)$, and the basis is carried along. Every relation is homogeneous of positive
length, so word length induces the grading, and the empty word spans the degree-zero part.
\end{proof}

Let $A_+(\Lambda)\subseteq A(\Lambda)$ be the span of the nonempty words. The proof of \cite[Corollary~3.8]{OAI247periodic}
shows, over any commutative ring, that the unital ideal $\mathcal J_\Lambda$ coincides with the two-sided ideal of the free
nonunital algebra $\Lambda\langle\Sigma\rangle_+$ generated by the same finite set of relations. Both ideals are spanned
by the products $xsy$, with $s$ a defining relation and $x,y$ possibly empty words. Consequently:

\begin{lemma}\label{lem:nonunital}
$A_+(\Lambda)$ is finitely presented as a nonunital $\Lambda$-algebra, by the generators $\Sigma$ and the relations
$\mathcal S$, and $A(\Lambda)=\Lambda\oplus A_+(\Lambda)$ is its unitization.
\end{lemma}

\subsection{Steinberg groups}

For a unital ring $R$ and $n\geq3$, the Steinberg group $\St_n(R)$ is generated by symbols $x_{ij}(a)$, where
$1\leq i\neq j\leq n$ and $a\in R$, subject to
\begin{align}
 x_{ij}(a+b)&=x_{ij}(a)x_{ij}(b),\label{eq:st-add}\\
 [x_{ij}(a),x_{kl}(b)]&=1\qquad (j\neq k,\ i\neq l),\label{eq:st-comm}\\
 [x_{ij}(a),x_{jk}(b)]&=x_{ik}(ab)\qquad (i,j,k\text{ pairwise distinct}),\label{eq:st-prod}
\end{align}
where $[g,h]=ghg^{-1}h^{-1}$. The assignment $x_{ij}(a)\mapsto e_{ij}(a)$, the elementary matrix with $a$ in position
$(i,j)$, defines a surjection $\pi_n:\St_n(R)\to\E_n(R)$ onto the elementary subgroup of $\GL_n(R)$. Its kernel is
$J_n(R)$. Write $\St(R)$, $\E(R)$ and $J(R)$ for the direct limits over $n$.

\subsection{The finite-rule hierarchy}

We use OpenAI's hierarchy of finite rule tables~\cite[Section~2]{OAI247periodic} without change. We recall what it
provides. Fix the absolute constants $c=10000$ and $k=10$ of \cite[Proposition~2.2]{OAI247periodic}, an integer
\begin{equation}\label{eq:D}
 D>10(c+1)(k+1),
\end{equation}
and an integer $h_0$, discussed in Remark~\ref{rem:parameters}. For $j\geq0$ put
\[
 h_j=h_0D^j,\qquad q_j=2^{h_j},\qquad L_j=2q_j+1,
\]
and let $\Sigma_j$ and $\mathcal S_j$ be the alphabet and rule table at parameter $h_j$. Write $A_j(\Lambda)$ for the
corresponding algebra.

\begin{proposition}[{\cite[Proposition~2.2]{OAI247periodic}}]\label{prop:hierarchy}
For every $j\geq0$ the following hold.
\begin{enumerate}[label=\textup{(H\arabic*)}]
\item\label{H:encoding} $|\Sigma_j|\leq2^{ch_j}$, and every rule of $\mathcal S_j$ has length at most $k$. There is a
  substitution $E_j:\Sigma_{j+1}\to\Sigma_j^{L_j}$ that carries every rule of $\mathcal S_{j+1}$ to a consequence of
  $\mathcal S_j$. It therefore induces a unital homomorphism $E_j:A_{j+1}(\Lambda)\to A_j(\Lambda)$. If $w$ is a nonzero
  word at level $j+1$, then $E_j(w)$ is a nonzero word at level $j$. Every level has a nonzero one-letter word.
\item\label{H:tokens} A designated set of letters marks \emph{cuts} immediately before them. There are a bound $B_j$ and a
  finite list $\mathcal T_j$ of \emph{tokens} such that every sufficiently long nonzero word at level $j$ has a unique
  decomposition $p\,t_1\cdots t_r\,s$ with $t_i\in\mathcal T_j$ and $|p|,|s|\leq B_j$. Here $p$ precedes the first cut, $s$
  starts at the last cut, and the $t_i$ are the intervals between consecutive cuts. For each $t\in\mathcal T_j$ there is a
  word $u_t$ over $\Sigma_{j+1}$ with $1\leq|u_t|\leq k$, $|t|=L_j|u_t|$ and $[t]=E_j([u_t])$.
\item\label{H:runs} Some letters are \emph{palette letters}, with a color in a set of $q_j$ colors. An \emph{eligible run}
  is a maximal string of consecutive palette letters with a common tuple of non-color fields and no control head. Every
  adjacent transposition inside an eligible run is an elementary replacement. Every eligible run in a nonzero word has
  pairwise distinct colors. Every sufficiently long nonzero word contains an eligible run of length at least $q_j/100$.
  Eligibility and the endpoints of maximal eligible runs depend only on the non-color fields.
\end{enumerate}
\end{proposition}

\begin{remark}[Independence of the field]\label{rem:char-free}
OpenAI state explicitly that a field is not part of the rule table: the same word relations are used over every
field~\cite[Section~2]{OAI247periodic}. Properties \ref{H:encoding}--\ref{H:runs} concern words, the replacement graph
and nonzero classes. By Lemma~\ref{lem:word-basis} these do not depend on the coefficient ring. We have checked that the
proofs in \cite[Section~2]{OAI247periodic} use no property of $\F_2$. Thus Proposition~\ref{prop:hierarchy} holds verbatim
over $\Z$ and over every field.
\end{remark}

\begin{remark}[Freedom in the parameters]\label{rem:parameters}
In \cite[Section~2.5]{OAI247periodic} the integer $D$ is fixed before the program fixed point. Afterwards $h_0$ is chosen
so large that, for every $h\geq h_0$, all field and layout bounds hold and every relevant computation finishes before the
tableau clock $T=2^{\lfloor h/5\rfloor}$ expires. Hence the conclusions of Proposition~\ref{prop:hierarchy} hold for every
integer $h_0$ exceeding a threshold that depends only on $D$ and the fixed program. Likewise $D$ may be any integer
satisfying~\eqref{eq:D}. Neither Theorem~A nor Theorem~B needs this freedom: we fix $D$ and $h_0$ once and for all, and
Section~\ref{sec:detection} handles every prime with the same table.
\end{remark}

%======================================================================
\section{The algebra in every positive characteristic}\label{sec:algebra}

We follow \cite[Section~3]{OAI247periodic}. Only one step depends on the characteristic: the cancellation of a symmetric
matrix sum inside a palette run. We replace it by a bound that depends on $p$ (Lemma~\ref{lem:symmetric}), and then check that the resulting
factor $p-1$ is absorbed by the hierarchy at a finite level.

Throughout this section $K$ is a field of characteristic $p>0$. The algebras $A_j(K)=K\otimes_{\Z}A_j(\Z)$ are those of
Section~\ref{sec:prelim}, and $E_j:A_{j+1}(K)\to A_j(K)$ is the homomorphism of~\ref{H:encoding}. We identify
$M_d(K)\otimes_KA_j(K)=M_d(A_j(K))$. A \emph{letter test} and a \emph{chunk test} of dimension $d$ (and label bound $b$)
at level $j$ are defined exactly as in \cite[Definition~3.1]{OAI247periodic}, with matrix weights in $M_d(K)$. Their
length-$n$ sums $S_{j,n}(M)$ and $C_{j,n}(\mathcal C,W)$ are as defined there. A test
\emph{eventually vanishes} if its length-$n$ sum is zero for all sufficiently large $n$.

\subsection{Symmetric sums in characteristic \texorpdfstring{$p$}{p}}

\begin{lemma}\label{lem:symmetric}
Let $X_1,\dots,X_m\in M_d(K)$, where $K$ has characteristic $p$, and put
\[
 \mathcal A_m(X_1,\dots,X_m)=\sum_{\sigma\in S_m}X_{\sigma(1)}\cdots X_{\sigma(m)}.
\]
If $m>(p-1)d^2$, then $\mathcal A_m(X_1,\dots,X_m)=0$.
\end{lemma}

\begin{proof}
Choose a basis $Y_1,\dots,Y_{d^2}$ of $M_d(K)$ and write $X_t=\sum_ic_{t,i}Y_i$. Since $\mathcal A_m$ is multilinear,
\[
 \mathcal A_m(X_1,\dots,X_m)=\sum_{(i_1,\dots,i_m)}\Bigl(\prod_{t=1}^mc_{t,i_t}\Bigr)
 \sum_{\sigma\in S_m}Y_{i_{\sigma(1)}}\cdots Y_{i_{\sigma(m)}} .
\]
Fix an index tuple $(i_1,\dots,i_m)$ and let $n_1,\dots,n_{d^2}$ be the multiplicities with which the basis indices
occur in it, so $\sum_in_i=m$. As $\sigma$ ranges over $S_m$, each distinct arrangement of these indices occurs exactly
$\prod_in_i!$ times. The inner sum therefore equals $\prod_in_i!$ times a sum of words in the $Y_i$. Since $m>(p-1)d^2$,
some $n_i\geq p$, so $p$ divides $\prod_in_i!$ and the inner sum vanishes in characteristic $p$.
\end{proof}

For $p=2$ this is the bound $m>d^2$ of \cite[Lemma~3.2]{OAI247periodic}. For $d\geq2$ the bound is far from sharp.
Zalesskii proved that $\mathcal A_{dp}$ vanishes identically on $M_d(K)$ in characteristic $p$, and that $\mathcal A_k$
does not for $k<dp$~\cite{Zalesskii1985}; we take the statement from~\cite{Valenti1998}. Grouping the terms of
$\mathcal A_{m+1}$ by their last factor writes it as a sum of terms $\mathcal A_m(\dots)X_t$, so $\mathcal A_m$ then
vanishes for every $m\geq dp$. Any bound that is finite for each $p$ and $d$ suffices below, and we keep
Lemma~\ref{lem:symmetric} because its proof is two lines. Its bound is sharp for $d=1$: for scalars equal
to $1$ the sum is $m!$, which is nonzero modulo $p$ when $m\leq p-1$. No signs or divisions occur. The quantity that
vanishes is a positive integer multiplicity, and this mechanism has no analogue in characteristic zero
(Section~\ref{sec:charzero}).

\subsection{Cancellation and its propagation}

\begin{lemma}[Direct cancellation in characteristic $p$]\label{lem:direct}
Fix $j,d,b$. If
\begin{equation}\label{eq:direct-bound}
 \frac{q_j}{100}>b\bigl((p-1)d^2+3\bigr),
\end{equation}
then every chunk test of dimension $d$ and label bound $b$ at level $j$ eventually vanishes over $K$.
\end{lemma}

\begin{proof}
The proof of \cite[Lemma~3.3]{OAI247periodic} applies with one change. It partitions the nonzero summands of a long
length-$n$ sum into free orbits of the symmetric group $S_m$. The group acts by permuting the $m$ complete chunks (the
\emph{packet}) inside the first eligible run of length at least $q_j/100$. Each orbit contributes
$U\,\mathcal A_m(X_1,\dots,X_m)\,V\otimes a$. That partition uses only~\ref{H:runs} and the replacement graph, and no
property of the field.

If the chosen run has length $r$ and its packet has $m$ chunks, then $r\leq mb+2b$, since at most two chunks cross its
endpoints. By~\eqref{eq:direct-bound}, $m\geq r/b-2>(p-1)d^2+1$. Hence every orbit contribution vanishes by
Lemma~\ref{lem:symmetric}, in place of \cite[Lemma~3.2]{OAI247periodic}.
\end{proof}

The next two results are \cite[Lemmas~3.4 and~3.5]{OAI247periodic}. Their proofs are linear algebra together
with~\ref{H:tokens}, and are valid over every field.

\begin{lemma}[{Reduction by the cuts~\cite[Lemma~3.4]{OAI247periodic}}]\label{lem:cut}
If every chunk test of dimension $d$ and label bound $k$ at level $j+1$ eventually vanishes, then every letter test of
dimension $d$ at level $j$ eventually vanishes.
\end{lemma}

\begin{lemma}[{A matrix automaton for chunks~\cite[Lemma~3.5]{OAI247periodic}}]\label{lem:automaton}
Let a chunk test over an alphabet of size $s$ have dimension $d$ and label bound $b$. There is a letter test of dimension
at most $b\,s^bd$ whose length-$n$ sum contains the chunk sum as a fixed $d\times d$ block, for every $n$.
\end{lemma}

\begin{proposition}[Eventual vanishing of every fixed test]\label{prop:all-tests}
Over every field $K$ of characteristic $p>0$, every letter test at level $0$ eventually vanishes, and so does every chunk
test at level $0$ with any fixed label bound.
\end{proposition}

\begin{proof}
We follow \cite[Proposition~3.6]{OAI247periodic}. Fix a starting dimension $d$, and define
$d_j=d\prod_{i=1}^{j-1}(k|\Sigma_i|^k)$. Lemmas~\ref{lem:cut} and~\ref{lem:automaton} reduce eventual vanishing of a
dimension-$d$ letter test at level $0$ to eventual vanishing of all dimension-$d_J$ chunk tests of label bound $k$ at
level $J$, for any $J\geq1$. As in the proof of \cite[Proposition~3.6]{OAI247periodic},
\[
 \log_2(d_j^2)\leq2\log_2d+2(j-1)\log_2k+\frac{2kc}{D-1}(h_j-h_1),
\]
and $\delta=1-2kc/(D-1)>0$ by~\eqref{eq:D}. Hence $h_j-\log_2(d_j^2)\to\infty$, and there is a finite $J$ with
\[
 2^{h_J}>400k(p-1)d_J^2\geq100k\bigl((p-1)d_J^2+3\bigr),
\]
where the second inequality uses $(p-1)d_J^2\geq1$. This is~\eqref{eq:direct-bound} with $b=k$. Lemma~\ref{lem:direct}
therefore applies at level $J$, and the finitely many reductions pull the vanishing back to level $0$, exactly as in
\cite{OAI247periodic}. The only difference from characteristic two is the factor $p-1$, which changes the finite stopping
level $J$ but not the presentation. A level-zero chunk test of any label bound is a block of a letter test, by
Lemma~\ref{lem:automaton}.
\end{proof}

\subsection{Matrix nilness and the proof of Theorem~A}

Let $\Sigma=\Sigma_0$ and $\mathcal S=\mathcal S_0$ be the level-zero alphabet and rule table, with $D$ and $h_0$ fixed
once and for all. This is the presentation of Theorem~A; it is written over $\Z$. For a field $K$ put $R_K=A_0(K)$ and
let $I_K=A_+(K)$ be its positive-degree ideal.

\begin{theorem}\label{thm:matrix-nil}
Let $K$ be a field of positive characteristic. Then $M_\ell(I_K)$ is nil for every $\ell\geq1$: every matrix with
entries in $I_K$ is nilpotent, with an index depending on the matrix.
\end{theorem}

\begin{proof}
This is the coefficient extraction of \cite[Theorem~3.7]{OAI247periodic}, run over an arbitrary field of characteristic
$p$. Write $T\in M_\ell(I_K)$ as $T=\sum_{\gamma}A_\gamma\otimes[v_\gamma]$, with $A_\gamma\in M_\ell(K)$ and
$|v_\gamma|\geq1$, and let $b=\max_\gamma|v_\gamma|$. Let $z$ be an indeterminate.

Apply Proposition~\ref{prop:all-tests} over the field $K(z)$, which has characteristic $p$, to the chunk test with labels
$v_\gamma$ and weights $zA_\gamma$. Its length-$n$ sum is the image of a polynomial $P_n(z)\in M_\ell(R_K)[z]$, as in the proof of
\cite[Theorem~3.7]{OAI247periodic}, and it vanishes for $n\geq N$. The map
$M_\ell(R_K)[z]\to M_\ell(K(z)\otimes_KR_K)$ is injective: expand coefficients in a $K$-basis of $R_K$ and use the
injectivity of $K[z]\to K(z)$. So $P_n(z)=0$ for $n\geq N$.

For $a\geq1$ we have $T^a=\sum_{n=a}^{ab}[z^a]P_n(z)$, since each factor has a label of length between $1$ and $b$.
Hence $T^a=0$ for $a\geq N$.
\end{proof}

The proof establishes cancellation afresh over every field, including $K(z)$. It does not rely on the general
principle, false for arbitrary nil algebras, that nilness survives extension of scalars.

\begin{proof}[Proof of Theorem~A]
\begin{itemize}[leftmargin=*]
\item \emph{Finite presentation.} By Lemma~\ref{lem:nonunital}, $I_K$ is finitely presented as a nonunital algebra by
  $\Sigma$ and $\mathcal S$, and $R_K=K\oplus I_K$ is its unitization, finitely presented by the same data.
\item \emph{Infinite dimension.} By~\ref{H:encoding}, the words $w_j=E_0E_1\cdots E_{j-1}(a_j)$, with $a_j$ a nonzero
  letter at level $j$, are nonzero and have strictly increasing lengths $L_0\cdots L_{j-1}$. By
  Lemma~\ref{lem:word-basis} they are nonzero over every field and linearly independent, so $I_K$ is
  infinite-dimensional.
\item \emph{Not nilpotent.} If $I_K^m=0$, every word of length at least $m$ would vanish, and $I_K$ would be spanned by the
  finitely many words of length less than $m$, a contradiction.
\item \emph{Nil and Jacobson radical.} By Theorem~\ref{thm:matrix-nil}, $M_\ell(I_K)$ is nil for every $\ell$. In
  particular $I_K$ is nil, hence Jacobson radical: if $i^N=0$ then $(1-i)^{-1}=\sum_{t<N}i^t$.
\item \emph{Algebraic unitization.} An element $a\cdot1+i$ of $R_K$ satisfies $(t-a)^N$ whenever $i^N=0$.
\end{itemize}
The presentation was fixed before $K$ was chosen, so the same presentation works for every field of positive
characteristic.
\end{proof}

%======================================================================
\section{Unsigned relations in characteristic zero}\label{sec:charzero}

\begin{proof}[Proof of Theorem~C]
Let $A$ be the nonunital algebra over a field $K$ of characteristic zero generated by a finite alphabet $\Sigma$, subject
to finitely many relations $u=v$ and $w=0$ with $u,v,w$ nonempty words. The proof of Lemma~\ref{lem:word-basis} does not
use equal lengths. Relations $u=v$ of different lengths still define a replacement graph, whose classes give a basis.
Hence the nonzero replacement classes of nonempty words form a $K$-basis of $A$.

Let $x=\sum_{a\in\Sigma}[a]$. Expanding,
\[
 x^n=\sum_{w\in\Sigma^n}[w]=\sum_{C}N_n(C)\,[C],
\]
where $C$ runs over the nonzero classes and $N_n(C)$ is the number of words of length $n$ lying in $C$. Each $N_n(C)$ is
a nonnegative integer, and it is positive exactly when $C$ contains a word of length $n$. In characteristic zero, positive
integers are nonzero in $K$. Therefore $x^n=0$ if and only if every word of length $n$ is zero in $A$.

Suppose that $A$ is nil. Then $x^N=0$ for some $N$, so every word of length $N$ vanishes. Every longer word has such a
word as a prefix, and a word with a zero prefix is zero. So every product of $N$ elements of $A$, being a combination of
words of length at least $N$, vanishes. Thus $A^N=0$.

For the presentation of Theorem~A, the words $w_j$ of the proof of Theorem~A are nonzero over every field, by
Lemma~\ref{lem:word-basis}. Their prefixes are nonzero as well, so nonzero words exist in every length. By the argument
above, $x^n\neq0$ for every $n$. In particular the algebra is not nil, and it is infinite-dimensional.
\end{proof}

\begin{remark}
\begin{enumerate}[label=(\arabic*),leftmargin=*]
\item \emph{What Theorem~C does not exclude.} It does not exclude finitely presented nil algebras that are not nilpotent
  in characteristic zero. It shows only that such an algebra cannot be presented by unsigned word relations alone.
  Relations with signs or other coefficients, such as $u+v=0$, binomials $u-\lambda v$, or relations with three or more
  terms, are outside its scope. Whether a finitely presented nil algebra that is not nilpotent exists over $\Q$ is the
  characteristic-zero case of Ufnarovskij's question~\cite[Question~8.1]{Smoktunowicz2008}, and it remains open.
\item \emph{Where the positive-characteristic mechanism lives.} In characteristic $p$, the coefficient of a word class in
  a power of a weighted sum is again a sum of positive integer multiplicities, and Lemma~\ref{lem:symmetric} arranges for
  these to be divisible by $p$. Theorem~C shows that no presentation by unsigned word relations can do this over $\Q$.
\item \emph{Semigroups versus algebras.} An algebra as in Theorem~C is the contracted semigroup algebra of the finitely
  presented semigroup with zero defined by the same relations. Nilness of the semigroup is much weaker than nilness of
  the algebra. Ivanov-Pogodaev and Kanel-Belov construct an infinite finitely presented nilsemigroup satisfying
  $x^9=0$~\cite{IvanovPogodaevKanelBelov2025}. By Theorem~C, its contracted semigroup algebra over any field of
  characteristic zero is not nil.
\item \emph{Priority.} The argument is the elementary observation that the sum of the generators has nonnegative
  integer coefficients on the word basis. We expect the statement to be known in the theory of semigroup
  algebras~\cite{Okninski1991}, and we do not claim it as new. We include it because it identifies exactly why the
  construction of Theorem~A is confined to positive characteristic.
\end{enumerate}
\end{remark}


%======================================================================
% Sections 5-8: the proof of Theorem B.
%======================================================================

\newcommand{\grpMat}{\operatorname{Mat}}
\newcommand{\grpdiag}{\operatorname{diag}}
\newcommand{\grpord}{\operatorname{ord}}
\newcommand{\grpsgn}{\operatorname{sgn}}
\newcommand{\grpbe}{\mathbf e}

%======================================================================
\section{Ordinary finite presentation over \texorpdfstring{$\F_p$}{F\_p}}\label{sec:fp}

Sections~\ref{sec:fp}--\ref{sec:detection} prove Theorem~B. The construction and the principal methods are those of
OpenAI's two manuscripts~\cite{OAI247periodic,OAI247rf}, which work over $\F_2$. At each step we say whether their
argument holds over $\F_p$ as written, in which case we state the result precisely and cite it, or what has to change, in
which case we give a complete proof. For $p=2$ every statement in these four sections is contained
in~\cite{OAI247periodic,OAI247rf}.

\paragraph{Conventions.}
Throughout Sections~\ref{sec:fp}--\ref{sec:detection}, $p$ is a fixed prime and $[g,h]=ghg^{-1}h^{-1}$. For a unital
ring $B$ we use the Steinberg groups $\St_n(B)$, with the relations~\eqref{eq:st-add}--\eqref{eq:st-prod} of
Section~\ref{sec:prelim}, the surjections $\pi_n\colon\St_n(B)\to\E_n(B)$, $x_{ij}(a)\mapsto e_{ij}(a)$, and their kernels
$J_n(B)$. Two root positions $(i,j)$ and $(k,l)$ are \emph{commuting positions} if $j\neq k$ and $i\neq l$; the
positions $(i,j)$ and $(j,i)$ are \emph{opposite}. Two consequences of the relations are used constantly:
$x_{ij}(a)^{-1}=x_{ij}(-a)$, and, since the right side of~\eqref{eq:st-prod} commutes with both factors, the reversed
commutator is $[x_{jk}(b),x_{ij}(a)]=x_{ik}(-ab)$. We write $E_{ij}$ for a matrix unit, $\grpbe_s$ for a standard column
and $\grpbe_s^*$ for the corresponding row, and $d_i(u)$ for the diagonal matrix with the unit $u$ in position $i$ and
$1$ elsewhere. The stable groups $\St(B)$, $\E(B)$ and $J(B)=\ker(\St(B)\to\E(B))$ are defined by the same presentation
with indices in $\{1,2,3,\dots\}$. Thus $\St(B)=\varinjlim_n\St_n(B)$, and every word and every derivation in it uses
finitely many indices. A word in roots, or in elementary and diagonal matrices, is \emph{supported} on a set $S$ of
indices if every index occurring in it lies in $S$.

Most of the argument concerns algebras of the following kind.

\begin{definition}\label{def:mn}
A unital $\F_p$-algebra $R$ is of \emph{matrix-nil type} if $R=\F_p\cdot1\oplus I$ for a two-sided ideal $I$ such that
every square matrix, of every finite size, with entries in $I$ is nilpotent.
\end{definition}

The algebras $R_p$ of Theorem~B are of matrix-nil type (Proposition~\ref{prop:det-input}), and so are their degree
truncations.

\begin{lemma}[Local structure; after {\cite[Lemmas~4.4 and~5.1]{OAI247periodic}}]\label{lem:fp-local}
Let $R$ be of matrix-nil type, and let $n\geq1$.
\begin{enumerate}[label=\textup{(\roman*)}]
\item A matrix in $M_n(R)$ is invertible if and only if its reduction modulo $I$ is invertible over $\F_p$. In
particular $R$ is a local ring with maximal ideal $I$.
\item Every tuple of columns of $R^n$ whose reductions are linearly independent over $\F_p$ is part of an invertible
matrix.
\item The group $\GL_n(R)$ is generated by the elementary matrices and the diagonal matrices $d_i(u)$, $u\in R^\times$.
\item Every element of $\GL_n(R)$ has finite order, and every element of $\ker\bigl(\GL_n(R)\to\GL_n(\F_p)\bigr)$ has
finite $p$-power order. In particular $\E_n(R)$ is periodic.
\end{enumerate}
\end{lemma}

\begin{proof}
(i) An invertible matrix has invertible reduction. Conversely, suppose the reduction $\bar A$ of $A$ is invertible, and
regard $C=\bar A^{-1}$ as a matrix over the copy of $\F_p$ in $R$. Then $AC=1+U$ and $CA=1+U'$ with $U,U'\in M_n(I)$.
If $U^m=0$, then $(1+U)^{-1}=\sum_{a<m}(-U)^a$, and similarly for $U'$. So $A$ has the right inverse $C(AC)^{-1}$ and
the left inverse $(CA)^{-1}C$, and is invertible. For $n=1$ this says that every element outside $I$ is a unit; the
elements of $I$ are nilpotent, hence not units.

(ii) Complete the reductions to a basis of $\F_p^n$, lift the added vectors, and apply (i).

(iii) Let $A$ be invertible. Some entry of its first column has nonzero reduction, and it is a unit by (i). The
matrix interchanging two coordinates $a,b$ equals $d_b(-1)\,W_{ab}$, where
\begin{equation}\label{eq:fp-W}
 W_{ab}=e_{ab}(1)\,e_{ba}(-1)\,e_{ab}(1)
\end{equation}
has the block $\left(\begin{smallmatrix}0&1\\-1&0\end{smallmatrix}\right)$ on the coordinates $a,b$, in this order.
Use it to move the unit $u$ to position $(1,1)$. Left multiplication by the matrices $e_{i1}(-a_{i1}u^{-1})$ and right
multiplication by the matrices $e_{1j}(-u^{-1}a_{1j})$ clear the rest of the first column and of the first row; the
orders of these products are valid in a noncommutative ring. The complementary block is invertible by (i), and
induction reduces $A$ to a diagonal matrix of units.

(iv) The reduction of $A$ lies in the finite group $\GL_n(\F_p)$, so $A^e=1+V$ for some $e\geq1$ and $V\in M_n(I)$; if
$A$ lies in the kernel, $e=1$. Choose $r$ with $V^r=0$ and $c$ with $p^c\geq r$. Since $1$ and $V$ commute and $p$
divides $\binom{p^c}{a}$ for $0<a<p^c$, we get $(1+V)^{p^c}=1+V^{p^c}=1$.
\end{proof}

For $p=2$ the interchange matrix is elementary, since $-1=1$. For odd $p$ it is not, because its reduction has
determinant $-1$. Most of the changes below come from this one fact.

\begin{lemma}[{\cite[Lemma~4.3]{OAI247periodic}}]\label{lem:fp-perfect}
For every unital ring $B$ and every $n\geq3$, the groups $\St_n(B)$ and $\E_n(B)$ are perfect, and both are infinite if
$B$ is infinite.
\end{lemma}

The proof in~\cite{OAI247periodic} is characteristic-free: $x_{ij}(a)=[x_{ih}(a),x_{hj}(1)]$ for $h\notin\{i,j\}$, and
distinct $a$ give distinct matrices $e_{ij}(a)$.

\subsection{The presentation}

Let $\mathcal A$ be a finite alphabet, $\mathcal A^*$ its free monoid, including the empty word $\varepsilon$, and
$T=\F_p\langle\mathcal A\rangle$ the free unital algebra, with basis $\mathcal A^*$. Every $f\in T$ has a unique
expression $f=\sum_w c_w(f)\,w$ with integers $c_w(f)\in\{0,1,\dots,p-1\}$, finitely many of them nonzero. Fix the
length-lexicographic order on $\mathcal A^*$. As in~\cite[Theorem~4.1]{OAI247periodic}, define a formal group word
$\xi_{ij}(w)$ in letters $g_{ij}(w)$, $|w|\leq2$, for every word $w$: put $\xi_{ij}(w)=g_{ij}(w)$ if $|w|\leq2$, and
for $|w|>2$ write $w=av$ with $a$ its first letter, let $h$ be the least index outside $\{i,j\}$, and put
$\xi_{ij}(w)=[\xi_{ih}(a),\xi_{hj}(v)]$.

\begin{theorem}[Ordinary finite presentation over $\F_p$]\label{thm:fp}
Let $R=T/\mathfrak a$, where $\mathfrak a=(f_1,\dots,f_s)$ is the two-sided ideal generated by finitely many elements,
and let $n\geq5$. Then $\St_n(R)$ has the ordinary finite presentation with generators $g_{ij}(w)$, for $i\neq j$ and
$|w|\leq2$, and relators
\begin{align}
 g_{ij}(w)^p&=1, && |w|\leq2,\label{eq:fp-order}\\
 [g_{ij}(u),g_{kl}(v)]&=1, && |u|,|v|\leq2,\ j\neq k,\ i\neq l,\label{eq:fp-comm}\\
 [g_{ij}(u),g_{jk}(v)]&=g_{ik}(uv), && |u|+|v|\leq2,\ i,j,k\text{ distinct},\label{eq:fp-prod}\\
 \prod_{w}\xi_{ij}(w)^{c_w(f_t)}&=1, && 1\leq t\leq s,\ i\neq j,\label{eq:fp-alg}
\end{align}
where the product in~\eqref{eq:fp-alg} is taken in the fixed order. The isomorphism sends $g_{ij}(w)$ to
$x_{ij}(\bar w)$, where $\bar w$ is the image of $w$ in $R$.
\end{theorem}

Finite presentability of Steinberg groups over finitely presented unital rings, for $n\geq4$, is due to Krsti\'c and
McCool~\cite[Theorem~3]{KrsticMcCool}; a finite presentation of $R$ as an $\F_p$-algebra becomes a finite ring
presentation once $p\cdot1=0$ is added. Theorem~\ref{thm:fp} is the characteristic-$p$ form of the explicit presentation
of~\cite[Theorem~4.1]{OAI247periodic}, which follows Krsti\'c and McCool's method of extending monomial root relations by
word length. Two things change: the square relators become $p$th powers, and the algebra relators carry
coefficients. The substantive point is the order of a new root. In characteristic two it follows from the roots being
involutions. Here it follows from the identity $C^p=[U^p,V]$, which requires that the new root $C=[U,V]$ commute with
its two factors. The induction is arranged so that these commutations are proved first.

\begin{lemma}[Three-factor identity; {\cite[Lemma~3.4]{OAI247rf}}, {\cite[Lemma~4.2]{OAI247periodic}}]
\label{lem:three-factor}
In any group, let $A=[X,Y]$ and $B=[Y,Z]$. If $[X,Z]=[A,X]=[B,Z]=[A,B]=1$, then $[A,Z]=[X,B]$, and this element
commutes with $A$, $Z$, $X$ and $B$.
\end{lemma}

\begin{proof}[Proof of Theorem~\ref{thm:fp}]
Omit~\eqref{eq:fp-alg} for now, and let $P_0$ be the group presented by~\eqref{eq:fp-order}--\eqref{eq:fp-prod}. We show
that $P_0\cong\St_n(T)$ and then impose the algebra relators. For $|w|\leq2$, write $y_{ij}(w)$ for the image of
$g_{ij}(w)$ in $P_0$.

\emph{The induction.} For an integer $N\geq2$ consider the assertions:
\begin{enumerate}[label=\textup{(I\arabic*)}]
\item elements $y_{ij}(w)$ are defined for all $|w|\leq N$, and $y_{ij}(w)=[y_{ih}(u),y_{hj}(v)]$ for every
factorization $w=uv$ into nonempty words and every $h\notin\{i,j\}$;
\item $y_{ij}(w)^p=1$ whenever $|w|\leq N$;
\item $[y_{ij}(u),y_{kl}(v)]=1$ whenever $|u|\leq N$, $|v|\leq N$, and $(i,j),(k,l)$ are commuting positions; there is
no condition on $|u|+|v|$, which may be as large as $2N$;
\item $[y_{ij}(u),y_{jk}(v)]=y_{ik}(uv)$ whenever $|u|+|v|\leq N$ and $i,j,k$ are distinct, including the cases
$u=\varepsilon$ and $v=\varepsilon$.
\end{enumerate}
For $N=2$ these are the defining relations. Assume them for $N-1\geq2$. We prove them for $N$ in four steps, in this
order.

\emph{Step 1: new roots.} This step is~\cite[proof of Theorem~4.1]{OAI247periodic} verbatim. For a word $w$ of length
$N$ and distinct $i,j$, every factorization $w=uv$ into nonempty words and every $h\notin\{i,j\}$ give a candidate
$[y_{ih}(u),y_{hj}(v)]$. To compare candidates, write $w=uvz$ with nonempty pieces, choose distinct indices $i,a,b,j$,
and apply Lemma~\ref{lem:three-factor} to $X=y_{ia}(u)$, $Y=y_{ab}(v)$, $Z=y_{bj}(z)$. By (I4) at $N-1$,
$A=y_{ib}(uv)$ and $B=y_{aj}(vz)$. Every root among $X,Y,Z,A,B$ has length less than $N$, and the four required
commutations are between commuting positions, so they hold by (I3) at $N-1$, even when the total length exceeds $N-1$.
The lemma gives $[y_{ib}(uv),y_{bj}(z)]=[y_{ia}(u),y_{aj}(vz)]$. This identifies any two candidates whose cut positions
and intermediate indices both differ. As $N-1\geq2$ cut positions and $n-2\geq3$ intermediate indices are available,
these identifications connect all candidates~\cite[proof of Theorem~4.1]{OAI247periodic}. Define $y_{ij}(w)$ as their
common value; by induction it is the element represented by $\xi_{ij}(w)$. This proves (I1) at $N$, and (I4) for
nonempty inputs of total length $N$. No order relation has been used.

\emph{Step 2: commuting pairs, before any order is asserted.} We prove (I3) for all pairs with $|u|,|v|\leq N$ by
induction on $|u|+|v|$. Pairs with both lengths below $N$ are covered by (I3) at $N-1$. Otherwise one of the two roots,
say $y_{ij}(u)$, has $|u|=N$. Let $y_{kl}(v)$ be the other, so that $j\neq k$, $i\neq l$ and $|v|\leq N$. Choose
$h\notin\{i,j,k,l\}$, which is possible as $n\geq5$, and a factorization $u=u_1u_2$ into nonempty words. By Step~1,
$y_{ij}(u)=[y_{ih}(u_1),y_{hj}(u_2)]$. The positions $(i,h)$ and $(h,j)$ are both commuting positions with $(k,l)$,
because $h\neq k$, $i\neq l$, $j\neq k$ and $h\neq l$. The totals $|u_1|+|v|$ and $|u_2|+|v|$ are smaller than
$|u|+|v|$, and all lengths involved are at most $N$. So $y_{kl}(v)$ commutes with both factors: by (I3) at $N-1$ if
$|v|<N$, and by the present induction if $|v|=N$. Hence it commutes with their commutator $y_{ij}(u)$.

\emph{Step 3: the order of a new root.} Let $|w|=N$ and write $C=y_{ij}(w)=[U,V]$ with $U=y_{ih}(u)$ and
$V=y_{hj}(v)$, for a factorization $w=uv$ into nonempty words and some $h\notin\{i,j\}$. The positions $(i,j)$ and
$(i,h)$ are commuting positions, and so are $(i,j)$ and $(h,j)$. All three words have length at most $N$, so Step~2
gives
\[
 [C,U]=[C,V]=1,
\]
although the total lengths $N+|u|$ and $N+|v|$ exceed $N$. In any group,
\[
 [U^k,V]=U\,[U^{k-1},V]\,U^{-1}\,[U,V].
\]
Since $C$ commutes with $U$, induction on $k$ gives $[U^k,V]=C^k$. Taking $k=p$ and using $U^p=1$, from (I2) at
$N-1$,
\[
 C^p=[U^p,V]=1.
\]
In characteristic two, \cite{OAI247periodic} uses instead that $U$ and $V$ are involutions, so that conjugation by $U$
inverts $C$. That argument has no odd analogue. The identity above replaces it and needs only the commutations of
Step~2, which were proved without any order relation.

\emph{Step 4: empty-word inputs.} This step is again~\cite[proof of Theorem~4.1]{OAI247periodic}. For a right empty
input, write $w=uv$ with $u,v$ nonempty, choose $h\notin\{i,j,k\}$, and apply Lemma~\ref{lem:three-factor} to
$X=y_{ih}(u)$, $Y=y_{hj}(v)$ and $Z=y_{jk}(\varepsilon)$. Here $A=y_{ij}(w)$ by Step~1 and $B=y_{hk}(v)$ by (I4) at
$N-1$. The required commutations are between commuting positions and words of length at most $N$, so they hold by
Step~2. The lemma gives $[y_{ij}(w),y_{jk}(\varepsilon)]=[y_{ih}(u),y_{hk}(v)]=y_{ik}(w)$. The left empty input is
symmetric, with $X=y_{ij}(\varepsilon)$, $Y=y_{jh}(u)$ and $Z=y_{hk}(v)$. This completes (I1)--(I4) at $N$.

\emph{From monomials to polynomials.} For $f\in T$ put
\[
 y_{ij}(f)=\prod_w y_{ij}(w)^{c_w(f)}.
\]
Roots in one position commute by (I3), so the order of the product is immaterial. Each $y_{ij}(w)$ has order dividing
$p$, so exponents add modulo $p$ and $y_{ij}(f+g)=y_{ij}(f)\,y_{ij}(g)$. The commuting relations for polynomials
follow from (I3). For distinct $i,j,k$, every commutator $[y_{ij}(u),y_{jk}(v)]=y_{ik}(uv)$ lies in position $(i,k)$,
which is a commuting position with both $(i,j)$ and $(j,k)$; so it commutes with every root in these two positions.
The identities $[XY,Z]=X[Y,Z]X^{-1}[X,Z]$ and $[X,YZ]=[X,Y]\,Y[X,Z]Y^{-1}$ therefore make the commutator
bimultiplicative on these roots, and
\[
 [y_{ij}(f),y_{jk}(g)]=\prod_{u,v}y_{ik}(uv)^{c_u(f)c_v(g)}=y_{ik}(fg),
\]
the last step by additivity, with coefficients reduced modulo $p$. Thus the elements $y_{ij}(f)$ satisfy the Steinberg
relations of $\St_n(T)$ and define $\Phi\colon\St_n(T)\to P_0$. Conversely, $g_{ij}(w)\mapsto x_{ij}(w)$ respects
\eqref{eq:fp-order}--\eqref{eq:fp-prod}, since $x_{ij}(w)^p=x_{ij}(pw)=1$, and defines $\Psi\colon P_0\to\St_n(T)$ with
$\Psi(y_{ij}(f))=x_{ij}(f)$. The two maps are inverse on generators, so $P_0\cong\St_n(T)$.

\emph{The algebra relators.} The left side of~\eqref{eq:fp-alg} represents $y_{ij}(f_t)$. So the quotient $P$ of $P_0$
by these relators is the quotient of $\St_n(T)$ by the roots $x_{ij}(f_t)$, in all positions. Write $\tilde x_{ij}(a)$
for the image of $x_{ij}(a)$ in $P$. For $c,d\in T$ and $h\notin\{i,j\}$,
\[
 \tilde x_{ij}(cf_t)=[\tilde x_{ih}(c),\tilde x_{hj}(f_t)]=1,\qquad
 \tilde x_{ij}(cf_td)=[\tilde x_{ih}(cf_t),\tilde x_{hj}(d)]=1,
\]
and additivity gives $\tilde x_{ij}(a)=1$ for every $a\in\mathfrak a$. The rest is as
in~\cite[proof of Theorem~4.1]{OAI247periodic}. The map $P\to\St_n(R)$ induced by $T\to R$ has the inverse
$x_{ij}(r)\mapsto\tilde x_{ij}(a)$, where $a$ is any lift of $r$; this is well defined because two lifts differ by an
element of $\mathfrak a$. There are finitely many generators and relators.
\end{proof}

%======================================================================
\section{Centrality and stability over \texorpdfstring{$\F_p$}{F\_p}}\label{sec:stability}

Throughout this section $R$ is of matrix-nil type (Definition~\ref{def:mn}). We prove that $J_n(R)$ is central in
$\St_n(R)$ for $n\geq6$ (Proposition~\ref{prop:centrality}), and that $\St_n(R)\to\St(R)$ is injective for $n\geq12$
(Theorem~\ref{thm:stability}).

The integral stability argument of~\cite[Section~7]{OAI247rf} rests on three finite-rank inputs, which
\cite[Proposition~7.1]{OAI247rf} imports from~\cite{OAI247periodic}: frame completion~\cite[Lemma~5.1]{OAI247periodic},
centrality~\cite[Proposition~5.4]{OAI247periodic}, and stability of $\GL_n/\E_n$~\cite[Lemma~5.5]{OAI247periodic}.
Their proofs use that a coordinate interchange is elementary in characteristic two: in the generation step
of~\cite[Lemma~5.1]{OAI247periodic}, in the elementary swap word $s$ of~\cite[Proposition~5.3]{OAI247periodic} on which
centrality rests, and in~\cite[Lemma~5.5]{OAI247periodic}, where elementary swap words are killed. We re-prove all
three inputs over $\F_p$: Lemma~\ref{lem:fp-local}, Propositions~\ref{prop:small-GL} and~\ref{prop:centrality}, and
Lemma~\ref{lem:GL-E}. The presentation of $\GL_n(R)$ uses true transpositions $s_{ab}=d_b(-1)w_{ab}$, which are
available there because diagonal generators are allowed. The quotient $\GL_n/\E_n$ uses the elementary signed
transpositions $w_{ab}$. These results replace~\cite[Proposition~7.1]{OAI247rf}.

For stability itself, \cite{OAI247periodic} proves rational stability through stabilizers of unordered frames. That
argument uses that the permutation block of an unordered frame stabilizer is elementary
\cite[proof of Lemma~5.9]{OAI247periodic}, and over an odd field it does not port
(Proposition~\ref{prop:unordered-fails}). We use instead the integral argument of~\cite[Section~7]{OAI247rf}, with
ordered frames and universal central extensions of pointwise stabilizers. Its one characteristic-two step is a face
comparison by a coordinate swap. We replace it with a signed elementary transposition that fixes the complementary block
exactly.

\paragraph{Signed transpositions.}
For distinct indices $a,b$, let $W_{ab}$ be the matrix~\eqref{eq:fp-W}, and let
\[
 w_{ab}=x_{ab}(1)\,x_{ba}(-1)\,x_{ab}(1)\in\St_n(R)
\]
be its standard lift. Thus $W_{ab}\grpbe_b=\grpbe_a$ and $W_{ab}\grpbe_a=-\grpbe_b$, and $W_{ab}$ fixes the other
standard columns. The product $s_{ab}=d_b(-1)W_{ab}$ is the transposition matrix of $a$ and $b$. For $p=2$ the sign
disappears and $W_{ab}=s_{ab}$. For odd $p$, $W_{ab}$ is elementary but $s_{ab}$ is not.

\subsection{A presentation of \texorpdfstring{$\GL_n(R)$}{GL\_n(R)} by small relations}

A \emph{frame} in $R^n$ is a tuple of columns whose reductions modulo $I$ are linearly independent over $\F_p$. Let
$X_n$ be the simplicial complex whose vertices are the columns with nonzero reduction and whose simplices are the finite
sets that form frames. Distinct columns are distinct vertices even when their reductions agree.

\begin{lemma}[{\cite[Lemma~5.2]{OAI247periodic}}]\label{lem:st-frames-sc}
For $n\geq3$ the complex $X_n$ is simply connected.
\end{lemma}

The proof in~\cite{OAI247periodic} adjoins labelled vectors one at a time and identifies each link with an independence
complex of projected vectors in a quotient of the residue vector space. It uses only linear algebra over the residue
field, and applies verbatim over $\F_p$.

\begin{proposition}[After {\cite[Proposition~5.3]{OAI247periodic}}]\label{prop:small-GL}
For $n\geq5$, the group $\GL_n(R)$ has the presentation with generators $e_{ij}(a)$, for $a\in R$ and $i\neq j$, and
$d_i(u)$, for $u\in R^\times$, subject to all words supported on sets of at most five coordinates whose matrix product
is the identity.
\end{proposition}

\begin{proof}
Let $\GL'_n$ be the presented group and $\pi\colon\GL'_n\to\GL_n(R)$ the natural map. It is surjective by
Lemma~\ref{lem:fp-local}(iii). We prove injectivity by induction on $n$. For $n=5$ every true relation is a relator.
Let $n>5$, and assume the result in rank $n-1$.

The one change from~\cite{OAI247periodic} is the swap word. Inside $\GL'_n$ we use the true transposition
\[
 s=d_2(-1)\,e_{12}(1)\,e_{21}(-1)\,e_{12}(1),
\]
a word supported on the coordinates $1,2$ whose matrix interchanges $\grpbe_1$ and $\grpbe_2$. In characteristic two
\cite{OAI247periodic} takes for $s$ an elementary word, which does not exist for odd $p$. Every relation among $s$ and
the other generators used below is a true matrix relation supported on at most four coordinates, hence a defining
relator. The argument never needs $s$ to be elementary.

\emph{The vertex stabilizer.} The stabilizer $P\leq\GL_n(R)$ of the column $\grpbe_1$ consists of the matrices
$\left(\begin{smallmatrix}1&b\\0&C\end{smallmatrix}\right)$ with $C\in\GL_{n-1}(R)$. Its complementary copy of
$\GL_{n-1}(R)$ lifts homomorphically to $\GL'_n$. By induction it is presented by true relations on at most five of the
coordinates $2,\dots,n$, and these are relators of $\GL'_n$. The row group generated by the $e_{1j}(a)$ also lifts,
because its additive and commuting relations involve at most three coordinates. The conjugation of the row group by
complementary generators is matrix conjugation, again by relations on at most four coordinates. Writing each element of
$P$ uniquely as a row transvection times $\grpdiag(1,C)$ gives a homomorphic section $P\to\GL'_n$ of $\pi$. Its image
$P'$ maps isomorphically onto $P$.

\emph{Generation.} We have $s^2=1$. Moreover $se_{i2}(a)s^{-1}=e_{i1}(a)$ for $i\geq3$,
$se_{12}(a)s^{-1}=e_{21}(a)$ and $sd_2(u)s^{-1}=d_1(u)$, all true relations on at most three coordinates. So $P'$ and
$s$ generate $\GL'_n$. Let $Q\leq P$ fix both $\grpbe_1$ and $\grpbe_2$. It is generated by the complementary
generators on the coordinates $3,\dots,n$ and the row generators $e_{ij}(a)$ with $i\in\{1,2\}$ and $j\geq3$. Let $Q'$ be
its image in $P'$. Conjugation by $s$ fixes the first kind and exchanges $e_{1j}(a)$ with $e_{2j}(a)$, so
$sQ's^{-1}=Q'$.

\emph{Transport on the frame complex.} The map $\GL'_n/P'\to X_n^{(0)}$, $hP'\mapsto\pi(h)\grpbe_1$, is well defined. If
$w$ is adjacent to $\pi(h)\grpbe_1$, frame completion (Lemma~\ref{lem:fp-local}(ii)) gives $q\in P$ with
$q\grpbe_2=\pi(h)^{-1}w$. Transport along the oriented edge to the coset $h\tilde q sP'$, where $\tilde q\in P'$ is the
section of $q$. Its vertex is $\pi(h)qs\grpbe_1=\pi(h)q\grpbe_2=w$. Here we use $s\grpbe_1=\grpbe_2$ exactly. The signed
$W_{12}$ would send $\grpbe_1$ to $-\grpbe_2$, which is a different vertex. The coset does not depend on $q$: if
$q\grpbe_2=q'\grpbe_2$, then $q'^{-1}q\in Q$ and $s^{-1}Q's=Q'\subseteq P'$. Nor does it depend on the representative
$h$: replacing $h$ by $hr$, with $r\in P'$, is offset by replacing $q$ by $\pi(r)^{-1}q$.

Transport reverses along an edge: from $h\tilde qs$ the reverse step may use $q=1$, and $h\tilde qs^2P'=hP'$.
Transport around a triangle is trivial. After a translation and frame completion, the triangle is
$(\grpbe_1,\grpbe_2,\grpbe_3)$. Let $t=d_3(-1)e_{23}(1)e_{32}(-1)e_{23}(1)$, the true transposition of $\grpbe_2$ and
$\grpbe_3$; its matrix lies in the complementary $\GL_{n-1}(R)$, so $t\in P'$. The successive representatives are $s$,
$sts$ and $ststs$, and $ststs=t$ is a true relation on three coordinates. So the endpoint coset is $P'$.

By Lemma~\ref{lem:st-frames-sc}, edge paths with the same endpoints differ by backtracks and triangles. Transport from
$P'$ therefore assigns one coset to each vertex. Every coset is reached, because $P'$ and $s$ generate $\GL'_n$:
right multiplication by $P'$ changes only the representative, and right multiplication by $s$ is an edge step with
$q=1$. The coset map is therefore a bijection. If $h\in\ker\pi$, then $hP'$ maps to $\grpbe_1$, so $h\in P'$; and $\pi$
is injective on $P'$. Hence $h=1$.
\end{proof}

\subsection{The central kernel and the quotient by elementary matrices}

\begin{proposition}[{\cite[Proposition~5.4]{OAI247periodic}} over $\F_p$]\label{prop:centrality}
For $n\geq6$ the kernel $J_n(R)$ is central in $\St_n(R)$.
\end{proposition}

\begin{proof}
This is the faithful-strip argument of~\cite[Proposition~5.4]{OAI247periodic}, a finite-rank form of Steinberg's proof of
stable centrality~\cite[Section~7]{Steinberg1967}. Given Proposition~\ref{prop:small-GL}, it holds over $\F_p$ as
written. We repeat it because the incoming-strip action carries a sign.

Let $e_{ij}(a)$ act on $\St_n(R)$ by conjugation by $x_{ij}(a)$, and let $d_k(u)$ act by the automorphism
$x_{ij}(c)\mapsto x_{ij}(u_icu_j^{-1})$, where $u_k=u$ and the other $u_i$ equal $1$; this automorphism respects the
relations, since $(u_iau_j^{-1})(u_jbu_k^{-1})=u_iabu_k^{-1}$. Let a true relation be supported on a set $S$ of at most
five coordinates, and choose $t\notin S$. The strips
\[
 U_t(v)=\prod_{i\in S}x_{it}(v_i),\qquad V_t(w)=\prod_{i\in S}x_{ti}(w_i)
\]
are abelian groups mapping injectively to matrices, because their parameters occupy distinct positions and no
products occur inside either strip. For $a,b\in S$, the relation~\eqref{eq:st-prod} and its reversed form give
\[
 x_{ab}(r)\,x_{bt}(v_b)\,x_{ab}(r)^{-1}=x_{at}(rv_b)\,x_{bt}(v_b),\qquad
 x_{ab}(r)\,x_{ta}(w_a)\,x_{ab}(r)^{-1}=x_{tb}(-w_ar)\,x_{ta}(w_a).
\]
So $e_{ab}(r)$ acts by $v\mapsto e_{ab}(r)v$ and by $w\mapsto w\,e_{ab}(r)^{-1}$; the minus sign in the incoming strip
is exactly the inverse. A diagonal $d_a(u)$ acts by $v_a\mapsto uv_a$ and $w_a\mapsto w_au^{-1}$. So a word with matrix
product $M$ supported on $S$ acts by
\[
 v\longmapsto Mv,\qquad w\longmapsto wM^{-1},
\]
and a true relation, with $M=1$, fixes every root having exactly one index in $S$. It fixes every root supported
outside $S$, and every root $x_{ij}(c)$ with $i,j\in S$, because $x_{ij}(c)=[x_{it}(c),x_{tj}(1)]$. By
Proposition~\ref{prop:small-GL} we obtain a homomorphism $\GL_n(R)\to\operatorname{Aut}\St_n(R)$ under which the image of
$\pi_n(z)$ acts as conjugation by $z$, for every $z\in\St_n(R)$. If $z\in J_n(R)$, then $\pi_n(z)=1$, and $z$ is
central.
\end{proof}

\begin{lemma}[{\cite[Lemma~5.5]{OAI247periodic}} over $\F_p$]\label{lem:GL-E}
The subgroup $\E_n(R)$ is normal in $\GL_n(R)$. For $n\geq5$, inclusion induces an isomorphism
\[
 \GL_n(R)/\E_n(R)\longrightarrow\GL_{n+1}(R)/\E_{n+1}(R).
\]
\end{lemma}

\begin{proof}
Elementary matrices normalize $\E_n(R)$, and $d_k(u)e_{ij}(a)d_k(u)^{-1}=e_{ij}(u_iau_j^{-1})$; so normality follows
from Lemma~\ref{lem:fp-local}(iii). By Proposition~\ref{prop:small-GL}, $\GL_n(R)/\E_n(R)$ is presented by the generators
$e_{ij}(a)$, $d_i(u)$, the true relations on at most five coordinates, and the relators $e_{ij}(a)$. Here the argument
of~\cite{OAI247periodic} kills elementary swap words; for odd $p$ there are none. We use the elementary signed
transposition instead. The true relation on two coordinates
\[
 W_{ij}\,d_i(u)\,W_{ij}^{-1}=d_j(u)
\]
holds for every unit $u$, even for $u$ not central, because the two signs it involves are central. Since $W_{ij}$ is a
product of elementary generators, the quotient identifies $d_i(u)$ with $d_j(u)$ for all $i,j$. So the quotient is
presented by one generator $d(u)$ for each unit $u$, with the relators obtained from true relations on at most five
coordinates by deleting the elementary letters and replacing each $d_i(u)$ by $d(u)$. Whether a word supported on $S$
is a true relation depends only on its matrix block on $S$, and every such word can be relabelled onto the coordinates
$1,\dots,5$. So this list of relators is the same for all $n\geq5$. Inclusion is the identity on the generators $d(u)$
and on their relators, which proves the assertion.
\end{proof}

\subsection{The unordered-frame argument does not port}

The rational stability argument of~\cite{OAI247periodic} uses, for $G=\E_m(R)$, $m\geq13$ and $1\leq l\leq4$, the
stabilizer $K_l\leq G$ of the \emph{unordered} frame $\{\grpbe_{m-l+1},\dots,\grpbe_m\}$. Its elements have the block form
$\left(\begin{smallmatrix}A&0\\C&P\end{smallmatrix}\right)$ with $A\in\GL_{m-l}(R)$ and $P$ a permutation matrix.
\cite[Lemma~5.9]{OAI247periodic} asserts that the projection onto the complementary block $A$ maps $K_l$ onto
$\E_{m-l}(R)$ with a locally finite kernel, and, for $l\geq2$, that the kernel contains a transposition of the frame. The transposition acts
by $-1$ on the rational orientation module and kills its coinvariants. Both assertions rest on permutation matrices
being elementary. The next proposition shows that both fail for odd $p$.

\begin{proposition}\label{prop:unordered-fails}
Let $p$ be odd, let $R$ be of matrix-nil type, and let $m\geq3$, $r=m-2$. Let $K\leq\E_m(R)$ be the stabilizer of the
unordered frame $\{\grpbe_{m-1},\grpbe_m\}$. Put
\[
 A_0=\grpdiag(-1,1,\dots,1)\in\GL_r(\F_p),\qquad P_2=\begin{pmatrix}0&1\\1&0\end{pmatrix},\qquad
 h=\grpdiag(A_0,P_2).
\]
\begin{enumerate}[label=\textup{(\roman*)}]
\item The matrix $h$ is a product of elementary matrices with entries in $\F_p$, so $h\in K$. Explicitly, with $W_{ab}$
as in~\eqref{eq:fp-W}, $W_{ab}(u)=e_{ab}(u)e_{ba}(-u^{-1})e_{ab}(u)$ and $H_{1m}=W_{1m}(-1)\,W_{1m}(1)^{-1}$,
\[
 H_{1m}=\grpdiag(-1,1,\dots,1,-1),\qquad h=H_{1m}\,W_{m-1,m}.
\]
On the three coordinates $1,m-1,m$, relabelled $1,2,3$, this reads $H_{13}W_{23}=\grpdiag(-1,P_2)$.
\item The complementary block $A_0$ of $h$ does not lie in $\E_r(R)$. So the complementary-block projection does not
map $K$ into $\E_r(R)$.
\item Every element of $K$ whose complementary block is the identity fixes $\grpbe_{m-1}$ and $\grpbe_m$. So the kernel of
the projection $K\to\GL_r(R)$ contains no transposition of the frame, and the orientation character is trivial on it.
\end{enumerate}
\end{proposition}

\begin{proof}
(i) On the coordinates $1,m$ in this order, $W_{1m}(u)$ has the block
$\left(\begin{smallmatrix}0&u\\-u^{-1}&0\end{smallmatrix}\right)$. So on these coordinates
\[
 W_{1m}(-1)\,W_{1m}(1)^{-1}=\begin{pmatrix}0&-1\\1&0\end{pmatrix}\begin{pmatrix}0&-1\\1&0\end{pmatrix}=
 \begin{pmatrix}-1&0\\0&-1\end{pmatrix},
\]
which gives $H_{1m}$. Next,
$W_{m-1,m}$ sends $\grpbe_m\mapsto\grpbe_{m-1}$ and $\grpbe_{m-1}\mapsto-\grpbe_m$ and fixes the other columns, and
$H_{1m}$ negates $\grpbe_1$ and $\grpbe_m$. Hence $H_{1m}W_{m-1,m}$ sends $\grpbe_1\mapsto-\grpbe_1$,
$\grpbe_{m-1}\mapsto\grpbe_m$ and $\grpbe_m\mapsto\grpbe_{m-1}$, and fixes the rest; that is, it equals $h$. The matrix
$h$ permutes $\grpbe_{m-1}$ and $\grpbe_m$, so it lies in $K$.

(ii) Reduction modulo $I$ maps $\E_r(R)$ into $\E_r(\F_p)=\SL_r(\F_p)$, and it fixes $A_0$, whose determinant is
$-1\neq1$.

(iii) Let $g=\left(\begin{smallmatrix}1&0\\C&P\end{smallmatrix}\right)\in K$. Its reduction lies in
$\E_m(\F_p)=\SL_m(\F_p)$ and has determinant $\det P$. So $\det P=1$, and $P$, a permutation matrix of size two, is the
identity.
\end{proof}

For $p=2$ the matrix $P_2$ is elementary and none of this occurs. Nor does the signed transposition repair the argument:
$W_{m-1,m}$ sends the frame vertex $\grpbe_{m-1}$ to $-\grpbe_m$, which is not a vertex of the frame, so it does not
stabilize the unordered frame. We therefore abandon unordered frames. In the ordered-frame argument below, stabilizers
are pointwise and contain no permutation block, and a signed transposition is needed only after one vertex has been
deleted.

\subsection{Universal central extensions of pointwise stabilizers}

For $r\geq6$ and $l\geq0$ let
\[
 \Pi_{r,l}=\left\{\begin{pmatrix}A&0\\C&1_l\end{pmatrix}:\ A\in\E_r(R),\ C\in\grpMat_{l\times r}(R)\right\}
 \leq\E_{r+l}(R).
\]
This is the semidirect product of the additive group of strips
$\left(\begin{smallmatrix}1&0\\C&1\end{smallmatrix}\right)$ with $\E_r(R)$, which acts by $C\mapsto CA^{-1}$. Let
$\widetilde\Pi_{r,l}$ be the semidirect product of the same strip group with $\St_r(R)$, acting through
$\pi_r$. The natural surjection $\widetilde\Pi_{r,l}\to\Pi_{r,l}$ has kernel $J_r(R)$. Recall that a
\emph{universal central extension} of a perfect group $G$ is a central extension $\widetilde G\to G$ that maps uniquely
over $G$ to every central extension of $G$. Its kernel is naturally $H_2(G;\Z)$, and this identification is natural for
the lifts of homomorphisms between universal central extensions; see~\cite[Section~5]{Milnor1971}
and~\cite[Section~III.5]{WeibelKbook}.

\begin{lemma}[{\cite[Lemma~7.3]{OAI247rf}} over $\F_p$]\label{lem:st-uce}
For $r\geq6$ and $l\geq0$, the map $\widetilde\Pi_{r,l}\to\Pi_{r,l}$ is a universal central extension.
Consequently the inclusion $A\mapsto\grpdiag(A,1_l)$ induces an isomorphism
\[
 H_2(\E_r(R);\Z)\xrightarrow{\ \sim\ }H_2(\Pi_{r,l};\Z),
\]
and both groups are naturally isomorphic to $J_r(R)$.
\end{lemma}

The proof of~\cite[Lemma~7.3]{OAI247rf} uses only the Steinberg relations, Lemma~\ref{lem:three-factor} and
centrality. It holds over $\F_p$ once centrality is available (Proposition~\ref{prop:centrality}). We give it in full,
because the signs of the strip action now matter.

\begin{proof}
The kernel $J_r(R)$ is central in $\St_r(R)$ by Proposition~\ref{prop:centrality}, and it acts trivially on the strips,
since the action is through $\pi_r$. So the extension is central. Both groups are perfect. An internal root is a
commutator $x_{ij}(a)=[x_{ih}(a),x_{hj}(1)]$, and a strip root has the same expression with an internal intermediate
index $h\leq r$.

Call the matrices $e_{ij}(a)\in\Pi_{r,l}$ with $j\leq r$, $i\leq r+l$ and $i\neq j$ the \emph{available roots};
they generate $\Pi_{r,l}$, and those with $i>r$ form the strip. Let $\rho\colon\mathcal X\to\Pi_{r,l}$ be a central
extension, and choose arbitrary lifts $y_{ij}(a)\in\mathcal X$ of the available roots.

\emph{Lifts in commuting positions commute.} Let $(i,j)$ and $(k,l')$ be available commuting positions. Choose $h\leq r$
outside $\{i,j,k,l'\}$. In matrices $e_{ij}(a)=[e_{ih}(a),e_{hj}(1)]$, and both factors commute with $e_{kl'}(b)$, since
$(i,h)$ and $(h,j)$ are commuting positions with $(k,l')$. In $\mathcal X$ the commutators of $y_{ih}(a)$ and $y_{hj}(1)$ with
$y_{kl'}(b)$ are therefore central. So conjugation by $y_{kl'}(b)$ multiplies both factors by central elements, and fixes
their commutator. That commutator differs from $y_{ij}(a)$ by a central element. Hence
\begin{equation}\label{eq:uce-commute}
 [y_{ij}(a),y_{kl'}(b)]=1.
\end{equation}
This holds for arbitrary lifts.

\emph{Corrected lifts.} Put $\hat y_{ij}(a)=[y_{ih}(a),y_{hj}(1)]$ for some $h\leq r$, $h\notin\{i,j\}$. For distinct
$i,h,j$ with $h\leq r$ and a fourth index $f\leq r$, apply Lemma~\ref{lem:three-factor} to $X=y_{ih}(a)$,
$Y=y_{hf}(b)$ and $Z=y_{fj}(1)$. The elements $[X,Y]$ and $[Y,Z]$ lift $e_{if}(ab)$ and $e_{hj}(b)$, and the four
hypotheses of the lemma are instances of~\eqref{eq:uce-commute}. Changing a lift by a central factor does not change a
commutator, so the lemma gives
\begin{equation}\label{eq:uce-product}
 [y_{ih}(a),y_{hj}(b)]=[y_{if}(ab),y_{fj}(1)].
\end{equation}
For $b=1$ this shows that $\hat y_{ij}(a)$ does not depend on $h$, and then~\eqref{eq:uce-product} reads
$[y_{ih}(a),y_{hj}(b)]=\hat y_{ij}(ab)$. Since $y_{ih}(a)y_{ih}(a')$ lifts $e_{ih}(a+a')$, and $y_{ih}(a)$ commutes with
$\hat y_{ij}(a')$ by~\eqref{eq:uce-commute},
\[
 \hat y_{ij}(a+a')=[y_{ih}(a)y_{ih}(a'),y_{hj}(1)]=y_{ih}(a)\,\hat y_{ij}(a')\,y_{ih}(a)^{-1}\,\hat y_{ij}(a)
 =\hat y_{ij}(a)\,\hat y_{ij}(a').
\]
Also $\hat y_{ij}(0)=1$, because a lift of the identity is central. The corrected lifts are themselves lifts, so
\eqref{eq:uce-commute} and~\eqref{eq:uce-product} apply to them as well. Hence they satisfy all the Steinberg relations
among available roots, and define homomorphisms $\St_r(R)\to\mathcal X$ and $\grpMat_{l\times r}(R)\to\mathcal X$.

\emph{The action.} For a strip root $e_{ij}(a)$, with $i>r$ and $j\leq r$, and an internal root $e_{jk}(b)$, the reversed
form of~\eqref{eq:uce-product} gives
\[
 \hat y_{jk}(b)\,\hat y_{ij}(a)\,\hat y_{jk}(b)^{-1}=\hat y_{ik}(-ab)\,\hat y_{ij}(a),
\]
which is the matrix rule $aE_{ij}\mapsto aE_{ij}\,e_{jk}(b)^{-1}=aE_{ij}-abE_{ik}$. All other internal roots commute with
the strip root, by~\eqref{eq:uce-commute}, and also in matrices. So the two homomorphisms respect the semidirect product
structure and define a homomorphism $\widetilde\Pi_{r,l}\to\mathcal X$ over $\Pi_{r,l}$. It is unique, since two
such homomorphisms differ by a homomorphism from the perfect group $\widetilde\Pi_{r,l}$ to the abelian group
$\ker\rho$. This proves universality.

For $l=0$ the lemma says that $\St_r(R)\to\E_r(R)$ is a universal central extension, so $J_r(R)=H_2(\E_r(R);\Z)$. The
inclusion $\E_r(R)\to\Pi_{r,l}$ lifts to $\St_r(R)\to\widetilde\Pi_{r,l}$, with zero strip, and this lift
is the identity on the kernels $J_r(R)$. Naturality of the identification of kernels with $H_2$ gives the isomorphism.
\end{proof}

\subsection{The complex of ordered frames}

An \emph{ordered frame} of length $q$ in $R^m$ is a tuple of $q$ columns whose reductions are linearly independent. Let
$\mathcal C_i(m)$ be the free abelian group on the ordered frames of length $i+1$, for $i\geq0$, with boundary the
alternating sum of the deletions. Let $\mathcal C_{-1}(m)=\Z$, generated by the empty frame, with augmentation sending
each frame of length one to $1$. Different orderings of the same columns are different generators; no relation
identifies them up to sign. Complexes of frames are classical in homological stability~\cite{vanderKallen1980}. The
finite complexes below are complexes of injective words on matroids, which Jonsson and Welker showed to be
shellable~\cite[Theorem~1.2 and Example~3.3]{JonssonWelker2009}. We follow the direct proof of~\cite[Lemma~7.4]{OAI247rf},
which is characteristic-free and gives the integral range we need.

\begin{lemma}[{\cite[Lemma~7.4]{OAI247rf}}]\label{lem:st-ordered}
The augmented complex $\mathcal C_\bullet(m)$ has $H_i=0$ for $-1\leq i\leq m-2$.
\end{lemma}

\begin{proof}
Let $V$ be a finite collection of labelled nonzero vectors spanning a vector space of dimension $m$ over $\F_p$. Labels
are kept distinct even when vectors agree. Let $\mathcal C_\bullet(V)$ be the augmented complex of ordered tuples of
distinct labels whose vectors are linearly independent. We prove $H_i(\mathcal C_\bullet(V))=0$ for $i\leq m-2$ by
induction on the number of labels. The empty collection has $m=0$, and the range is empty.

Delete a label $v$, put $V'=V\setminus\{v\}$, and consider the relative complex
$Q_\bullet=\mathcal C_\bullet(V)/\mathcal C_\bullet(V')$, spanned by the independent tuples containing $v$. Filter it by
the number $a$ of labels preceding $v$. Deleting $v$ gives zero in $Q_\bullet$, deleting an earlier label lowers $a$, and
deleting a later label preserves $a$. So the associated graded complex at $a$ is a direct sum, over independent ordered
prefixes $(u_1,\dots,u_a,v)$, of complexes in which only the suffix varies. Project the remaining labelled vectors to
the quotient by $\langle u_1,\dots,u_a,v\rangle$ and discard the zero projections. A suffix extends the prefix to an
independent tuple exactly when its projections are independent. So the summand is the augmented complex of the
projected collection, shifted up by $a+1$ and with boundary multiplied by $(-1)^{a+1}$. The projected collection spans
dimension $m-a-1$ and has fewer labels, so by induction the summand has no homology in degrees
$\leq(m-a-1)-2+(a+1)=m-2$. When $m-a-1=0$ its only homology is in degree $-1+(a+1)=m-1$, also outside this range. As the
filtration is finite, $H_i(Q_\bullet)=0$ for $i\leq m-2$.

If $V'$ still spans dimension $m$, the long exact sequence of the pair and the induction hypothesis for $V'$ finish the
proof. Otherwise $v$ is independent of the span of $V'$. Then prepending $v$, $t\mapsto(v,t)$, with the empty tuple sent
to $(v)$, is a chain homotopy from the inclusion $\mathcal C_\bullet(V')\to\mathcal C_\bullet(V)$ to zero; indeed
$\partial(v,t)=t-(v,\partial t)$. So the inclusion is zero on homology. It is also surjective in degrees $\leq m-2$, by the
vanishing for $Q_\bullet$. Hence $H_i(\mathcal C_\bullet(V))=0$ for $i\leq m-2$.

Finally, a cycle or chain of $\mathcal C_\bullet(m)$ uses finitely many columns. Adjoin the standard columns, and take
the reductions as labelled vectors. The finite collection spans $\F_p^m$, and its complex is the subcomplex of frames
using only these columns. The finite case therefore applies.
\end{proof}

\subsection{Stability of \texorpdfstring{$J_n$}{J\_n} and injectivity at fixed rank}

\begin{theorem}[{\cite[Theorem~7.2]{OAI247rf}} over $\F_p$]\label{thm:stability}
Let $R$ be of matrix-nil type. For $m\geq13$ the stabilization map $J_{m-1}(R)\to J_m(R)$ is an isomorphism.
Consequently $\St_n(R)\to\St(R)$ is injective for every $n\geq12$.
\end{theorem}

For the characteristic-two ring of~\cite{OAI247rf}, its Section~7 notes that the injectivity also follows from the
stability theorem of Suslin and Tulenbaev~\cite{SuslinTulenbaev1976}. We follow instead the direct integral argument
of~\cite[Section~7]{OAI247rf}, which makes every natural map explicit. Only the face comparison below changes.

\begin{proof}
Fix $m\geq13$ and write $\E_m=\E_m(R)$.

\emph{Transitivity and stabilizers.} Let $1\leq l\leq5$ and $r=m-l\geq8$. Given an ordered frame of length $l$,
Lemma~\ref{lem:fp-local}(ii) completes it to $g\in\GL_m(R)$ whose last $l$ columns are the frame. By
Lemma~\ref{lem:GL-E}, applied in the ranks $r,\dots,m-1$, the class of $g$ in $\GL_m(R)/\E_m$ contains
$\grpdiag(A,1_l)$ for some $A\in\GL_r(R)$. So $g\,\grpdiag(A^{-1},1_l)\in\E_m$, and it has the same last $l$ columns.
Thus $\E_m$ acts transitively on ordered frames of length $l$. The pointwise stabilizer of $(\grpbe_{r+1},\dots,\grpbe_m)$
in $\GL_m(R)$ consists of the matrices $\left(\begin{smallmatrix}A&0\\C&1_l\end{smallmatrix}\right)$ with
$A\in\GL_r(R)$. Such a matrix equals $\left(\begin{smallmatrix}1&0\\CA^{-1}&1\end{smallmatrix}\right)\grpdiag(A,1_l)$,
whose first factor is elementary. So it lies in $\E_m$ if and only if $A\in\E_r(R)$, by the injectivity in
Lemma~\ref{lem:GL-E}. The stabilizer in $\E_m$ is therefore exactly $\Pi_{r,l}$. It contains no permutation of the
frame; compare Proposition~\ref{prop:unordered-fails}. Write $\Pi_l=\Pi_{m-l,l}$.

\emph{The spectral sequence.} Let $P_\bullet\to\Z$ be a free resolution by right $\Z[\E_m]$-modules, and form the
double complex $D_{i,q}=P_q\otimes_{\Z[\E_m]}\mathcal C_i(m)$, for $i,q\geq0$. Augmented by
$P_q\otimes_{\Z[\E_m]}\Z$, each row is exact in degrees $\leq m-2$, by Lemma~\ref{lem:st-ordered} and freeness of $P_q$.
Hence augmentation identifies the homology of the total complex with $H_j(\E_m;\Z)$ for $j\leq m-3$, in particular for
$j=2$; see~\cite[Chapter~VII]{Brown1982}, or the explicit elimination in~\cite[proof of Proposition~5.10]{OAI247periodic},
which uses only row exactness and freeness. Filtering by columns gives a spectral sequence $E^1_{i,q}$ converging to this
homology. For $0\leq i\leq4$, $\mathcal C_i(m)$ is the permutation module on one orbit of ordered frames of length $i+1$,
with stabilizer $\Pi_{i+1}$ of the standard frame $(\grpbe_{m-i},\dots,\grpbe_m)$. By Shapiro's lemma
\[
 E^1_{i,q}=H_q(\Pi_{i+1};\Z)\qquad(0\leq i\leq4).
\]
The differential $d^1$ is the alternating sum of the face maps. The $k$th face map is induced by
$\Pi_{i+1}\to\Pi_i$, $x\mapsto g^{-1}xg$, where $g\in\E_m$ carries the standard frame of length $i$ to
the $k$th face of the standard frame of length $i+1$. Different choices of $g$ differ by elements of $\Pi_i$,
which act trivially on its homology.

\emph{Rows $q=0$ and $q=1$.} In row $0$ every entry is $\Z$ and every face map is the identity, so
$d^1\colon E^1_{i,0}\to E^1_{i-1,0}$ is zero for odd $i$ and the identity for even $i\geq2$. Hence
$E^2_{1,0}=E^2_{2,0}=E^2_{3,0}=0$; the last uses the identity map from column four, that is, transitivity on frames of
length five. In row $1$, $E^1_{i,1}=H_1(\Pi_{i+1};\Z)=0$ for $i\leq4$, because these groups are perfect
(Lemma~\ref{lem:st-uce}). So $E^2_{1,1}=E^2_{2,1}=0$.

\emph{The face comparison (the step that changes).} The two face maps $H_2(\Pi_2;\Z)\to H_2(\Pi_1;\Z)$
start from the standard frame $(\grpbe_{m-1},\grpbe_m)$. Deleting its first entry leaves $(\grpbe_m)$, the standard
frame of length one, and this face map is induced by the inclusion $\Pi_2\subseteq\Pi_1$. Deleting its
second entry leaves $(\grpbe_{m-1})$. In~\cite{OAI247rf} this column is returned to $\grpbe_m$ by interchanging the last
two coordinates, which is elementary only in characteristic two. We use instead the elementary signed transposition
\[
 W=W_{m,m-1}=e_{m,m-1}(1)\,e_{m-1,m}(-1)\,e_{m,m-1}(1)\in\E_m(\F_p)\subseteq\E_m,
\]
whose standard lift is $w_{m,m-1}$. It satisfies $W\grpbe_{m-1}=\grpbe_m$ exactly and $W\grpbe_m=-\grpbe_{m-1}$, and it fixes
$\grpbe_1,\dots,\grpbe_{m-2}$. Taking $g=W^{-1}$, the second face map is induced by $x\mapsto WxW^{-1}$. The sign falls
on $\grpbe_m$, the column that was deleted, and plays no role. Since $W=\grpdiag(1_{m-2},W')$ is block diagonal,
\[
 W\,\grpdiag(A,1_2)\,W^{-1}=\grpdiag(A,1_2)\qquad(A\in\E_{m-2}(R)).
\]
So the two face maps restrict to the same homomorphism on the complementary block $\E_{m-2}(R)\leq\Pi_2$. In the
Steinberg group the same holds exactly: the three factors of $w_{m,m-1}$ commute with every root $x_{ij}(a)$,
$i,j\leq m-2$. By Lemma~\ref{lem:st-uce}, with $r=m-2$ and $l=2$, $H_2(\E_{m-2}(R);\Z)\to H_2(\Pi_2;\Z)$ is an
isomorphism. So the two face maps agree on all of $H_2(\Pi_2;\Z)$, $d^1=0$ on $E^1_{1,2}$, and
\[
 E^2_{0,2}=H_2(\Pi_1;\Z).
\]
Proposition~\ref{prop:unordered-fails} shows that a signed transposition cannot stand in for a transposition when both
vertices of a frame must be kept. Here only one vertex survives the deletion, and the signed transposition suffices.

\emph{Conclusion.} No differential leaves column $0$. The only differentials that could reach position $(0,2)$ come from
$E^2_{2,1}=0$ and from $E^3_{3,0}$, a subquotient of $E^2_{3,0}=0$. The other positions of total degree two have
$E^2_{1,1}=E^2_{2,0}=0$. Hence the edge map $H_2(\Pi_1;\Z)\to H_2(\E_m;\Z)$ is an isomorphism. It is induced by
the inclusion $\Pi_1\subseteq\E_m$, because column $0$ corresponds to the vertex $\grpbe_m$ and the augmentation
$P_\bullet\otimes_{\Z[\Pi_1]}\Z\to P_\bullet\otimes_{\Z[\E_m]}\Z$ is the map induced by inclusion. By
Lemma~\ref{lem:st-uce} with $r=m-1$ and $l=1$, the inclusion $\E_{m-1}(R)\to\Pi_1$ is an isomorphism on $H_2$.
The composite $H_2(\E_{m-1}(R);\Z)\to H_2(\E_m(R);\Z)$ is induced by the standard inclusion and is therefore an
isomorphism. The stabilization $\St_{m-1}(R)\to\St_m(R)$ lifts this inclusion between universal central extensions
(Lemma~\ref{lem:st-uce} with $l=0$). By naturality, $J_{m-1}(R)\to J_m(R)$ is an isomorphism.

Finally, let $n\geq12$, and let $z\in\St_n(R)$ map to $1$ in $\St(R)$. Since $\St(R)$ is a direct limit, $z$ already
maps to $1$ in some $\St_{n'}(R)$. The matrix groups embed in one another, so $\pi_n(z)=1$ and $z\in J_n(R)$. The
composite $J_n(R)\to J_{n'}(R)$ is a composite of isomorphisms, so $z=1$.
\end{proof}

%======================================================================
\section{Relative \texorpdfstring{$p$}{p}-primary torsion}\label{sec:torsion}

In this section we prove the following theorem.

\begin{theorem}[Relative $p$-primary torsion]\label{thm:relative-torsion}
Let $R=\bigoplus_{d\geq0}R_d$ be a graded unital $\F_p$-algebra with $R_0=\F_p$. Every element of
\[
 J(R)\cap\ker\bigl(\St(R)\longrightarrow\St(\F_p)\bigr)
\]
has finite $p$-power order.
\end{theorem}

For $p=2$ this is~\cite[Proposition~4.5]{OAI247rf}; the underlying comparison is the stable-torsion argument
of~\cite[Lemmas~6.4--6.7 and Proposition~6.8]{OAI247periodic}. The grading lets one substitute a numerical $b\times b$
matrix for a central variable. These substitutions are elementary versions of Weibel's operations on the $K$-theory of
graded rings~\cite[Section~1]{Weibel1987}. Evaluation at a long cyclic shift is trivial, and evaluation at a cyclic shift
of suitable size is the $b$th power of the element. Nothing about $R$ beyond the grading is used: neither matrix-nilness
nor the hierarchy. Three points change for odd $p$:
\begin{itemize}[leftmargin=*]
\item $b$ must be a power of $p$, so that $(C_b-1)^b=0$;
\item the basis change that triangularizes the cyclic shift must be corrected to determinant one, which is automatic over
$\F_2$;
\item coordinate permutations are realized with a compensating sign on an unused coordinate. The characteristic-two lemma
that lifts a transposition is false for odd $p$, and Lemma~\ref{lem:tor-relabel} replaces it.
\end{itemize}
The conclusion concerns only central elements that die over $\F_p$. Nothing about the torsion of $J(\F_p)$ is used or
claimed.

\subsection{Strips, centrality and coordinate changes}

\begin{lemma}[Faithful strips; {\cite[Lemma~3.1]{OAI247rf}}]\label{lem:tor-strips}
Let $B$ be a unital ring, and let $S,T$ be disjoint finite sets of indices. The subgroup of $\St(B)$ generated by the
roots $x_{ij}(a)$, $i\in S$, $j\in T$, is abelian, and $u\mapsto\pi(u)-1$ maps it isomorphically onto the additive group
of $S\times T$ matrices over $B$. A word supported on $S$, with matrix $M$ on $S$, normalizes this subgroup and acts on it
by $X\mapsto MX$. On the reverse strip $T\times S$ it acts by $X\mapsto XM^{-1}$.
\end{lemma}

The proof in~\cite{OAI247rf} is characteristic-free; the inverse in the reverse action comes from the sign of the
reversed commutator, as in the proof of Proposition~\ref{prop:centrality}.

\begin{lemma}[Stable centrality; {\cite[Lemma~3.2]{OAI247rf}}]\label{lem:tor-center}
For every unital ring $B$, the subgroup $J(B)$ is central in $\St(B)$.
\end{lemma}

\begin{proof}
Let $z\in J(B)$ be represented by a word supported on a finite set $S$, which we enlarge to contain the indices of a given
root $x_{ij}(a)$. Choose $h\notin S$. By Lemma~\ref{lem:tor-strips}, $z$ acts trivially on the strips $S\times\{h\}$ and
$\{h\}\times S$, since its matrix is the identity. So $z$ commutes with $x_{ij}(a)=[x_{ih}(a),x_{hj}(1)]$.
\end{proof}

\begin{lemma}\label{lem:tor-SL}
For $b\geq1$, every matrix in $\SL_b(\F_p)$ is a product of elementary matrices $e_{ij}(\lambda)$, $\lambda\in\F_p$.
\end{lemma}

\begin{proof}
This is standard; we record the factors used. Let $Q\in\SL_b(\F_p)$ with $b\geq2$. If $Q_{11}=0$, adding a suitable row
to the first makes it nonzero. Elementary row and column operations then clear the rest of the first column and row,
and induction reduces $Q$ to a diagonal matrix $\grpdiag(u_1,\dots,u_b)$ with $\prod u_i=1$. With $v_k=u_1\cdots u_k$, this
diagonal matrix is the product over $k<b$ of the matrices with $v_k$ and $v_k^{-1}$ in positions $k$ and $k+1$. Each of
these is $W_{k,k+1}(v_k)\,W_{k,k+1}(1)^{-1}$, where $W_{ab}(u)=e_{ab}(u)e_{ba}(-u^{-1})e_{ab}(u)$ has the block
$\left(\begin{smallmatrix}0&u\\-u^{-1}&0\end{smallmatrix}\right)$.
\end{proof}

In characteristic two, \cite[Lemma~6.3]{OAI247periodic} and~\cite[Lemma~3.3]{OAI247rf} lift the transposition of
coordinates $a,b$ to $x_{ab}(1)x_{ba}(1)x_{ab}(1)$. For odd $p$ that element does not induce the transposition. Its
matrix on the coordinates $a,b$ is $\left(\begin{smallmatrix}2&3\\1&2\end{smallmatrix}\right)$, which is
$\left(\begin{smallmatrix}2&0\\1&2\end{smallmatrix}\right)$ over $\F_3$. No element of $\E_n(\F_p)$ has a transposition
as its matrix, since its determinant would be $-1$. The characteristic-two lemma is used in four places: to show that the
evaluation $W_z(T)$ below does not depend on the coordinates chosen, to identify the coordinate copies $\iota_s(z)$ with
$z$, in the first step of relocation (Lemma~\ref{lem:det-relocation}), and to identify the offset words in
Proposition~\ref{prop:det-power}. We replace it by the following lemma, in which the sign is compensated on an unused
coordinate.

\begin{lemma}[Relabelling coordinates over $\F_p$; replaces {\cite[Lemma~3.3]{OAI247rf}} and
{\cite[Lemma~6.3]{OAI247periodic}}]\label{lem:tor-relabel}
Let $B$ be a unital $\F_p$-algebra, $S$ a finite set of indices, and $f\colon S\to\{1,2,3,\dots\}$ an injective map.
There is an element $g\in\St(B)$, a product of roots with entries in $\F_p$, such that
\[
 g\,x_{ij}(a)\,g^{-1}=x_{f(i)f(j)}(a)\qquad(i,j\in S,\ i\neq j,\ a\in B).
\]
Consequently, if $z\in J(B)$ is represented by a word supported on $S$, the word obtained by applying $f$ to every index
also represents $z$.
\end{lemma}

\begin{proof}
Put $S'=S\cup f(S)$. The sets $S'\setminus S$ and $S'\setminus f(S)$ have the same size, so $f$ extends to a
permutation $\sigma$ of $S'$. Choose distinct indices $c,h\notin S'$. Let $g_0$ be the numerical matrix with
$g_0\grpbe_k=\grpbe_{\sigma(k)}$ for $k\in S'$, $g_0\grpbe_c=\grpsgn(\sigma)\,\grpbe_c$, and $g_0\grpbe_k=\grpbe_k$
otherwise. It has determinant one and is supported on $S''=S'\cup\{c\}$. By Lemma~\ref{lem:tor-SL} it is a product of
elementary matrices $e_{kl}(\lambda)$ with $k,l\in S''$ and $\lambda\in\F_p$. Let $g\in\St(B)$ be the product of the
corresponding roots $x_{kl}(\lambda)$. It is supported on $S''$, and $h\notin S''$.

By Lemma~\ref{lem:tor-strips}, $g$ acts on the strip $S''\times\{h\}$ by left multiplication by $g_0$. For $i\in S$ the
root $x_{ih}(a)$ is the column $a\grpbe_i$, and $g_0(a\grpbe_i)=a\grpbe_{\sigma(i)}$ carries no sign, since $i\neq c$.
So $g\,x_{ih}(a)\,g^{-1}=x_{\sigma(i)h}(a)$. On the reverse strip, $g$ acts by right multiplication by $g_0^{-1}$; for
$j\in S$ the row $\grpbe_j^*$ goes to $\grpbe_j^*g_0^{-1}=\grpbe_{\sigma(j)}^*$, so
$g\,x_{hj}(1)\,g^{-1}=x_{h\sigma(j)}(1)$. Using $x_{ij}(a)=[x_{ih}(a),x_{hj}(1)]$ for distinct $i,j\in S$,
\[
 g\,x_{ij}(a)\,g^{-1}=[x_{\sigma(i)h}(a),x_{h\sigma(j)}(1)]=x_{\sigma(i)\sigma(j)}(a).
\]
This is an exact identity in $\St(B)$, not merely an identity of matrices. The sign compensation sits on $c$, which is
not an endpoint of any root supported on $S$. For the last assertion, $gzg^{-1}=z$ by Lemma~\ref{lem:tor-center}, and
$gzg^{-1}$ is represented by the relabelled word.
\end{proof}

An equivalent route uses only signed transpositions. The product $w_{bc}^2\,w_{ab}\,w_{bc}$ has the matrix of the
unsigned $3$-cycle $\grpbe_a\mapsto\grpbe_b\mapsto\grpbe_c\mapsto\grpbe_a$, and by Lemma~\ref{lem:tor-strips} it acts on
roots with exactly that permutation of indices. Three-cycles generate the alternating group, and any injective relabelling
of a finite set extends to an even permutation of a larger finite set.

\subsection{Block roots and triangular evaluation}

Fix a finite set of \emph{outer} indices and $b\geq1$, and label coordinates by pairs $(i,s)$, with $i$ outer and $s$
an \emph{inner} index in $\{0,\dots,b-1\}$; in the stable group we fix an injection of these pairs into the index set.
For a unital $\F_p$-algebra $B$, distinct outer indices $i,j$ and $A=(a_{st})\in\grpMat_b(B)$, the \emph{block root} is
\[
 X_{ij}(A)=\prod_{s,t}x_{(i,s),(j,t)}(a_{st}).
\]
Its factors commute, since all sources have outer index $i$ and all targets the different outer index $j$.

\begin{lemma}[Block roots; after {\cite[Lemma~4.1]{OAI247rf}} and {\cite[Lemma~6.4]{OAI247periodic}}]\label{lem:tor-block}
\begin{enumerate}[label=\textup{(\roman*)}]
\item The block roots satisfy the Steinberg relations over $\grpMat_b(B)$, and so define a homomorphism from the
Steinberg group over $\grpMat_b(B)$, on the outer indices, to $\St(B)$.
\item For every $Q\in\SL_b(\F_p)$ there is $q\in\St(B)$, a product of roots with entries in $\F_p$ supported inside the
outer blocks, such that
\[
 q\,X_{ij}(A)\,q^{-1}=X_{ij}(QAQ^{-1})\qquad\text{for all distinct outer }i,j\text{ and all }A.
\]
Its matrix is block diagonal, with a copy of $Q$ in each outer block.
\end{enumerate}
\end{lemma}

\begin{proof}
Part (i) is~\cite[Lemma~4.1]{OAI247rf}; its proof is characteristic-free. In the consecutive commutator, every resulting
root commutes with both input blocks, so the commutator expands into the matrix product.

For (ii), the case $b=1$ is trivial. In characteristic two \cite{OAI247rf} allows every $Q\in\GL_b(\F_2)$. Over $\F_p$
we need determinant one, and only that case is used below. Write $Q=e_1\cdots e_k$ with elementary factors over $\F_p$
(Lemma~\ref{lem:tor-SL}). For an elementary $e=e_{st}(\lambda)$, let $\hat e$ be the product over all outer indices $o$
of the roots $x_{(o,s),(o,t)}(\lambda)$; these commute, having disjoint supports. By Lemma~\ref{lem:tor-strips}, the
factor in block $i$ acts on the strip from block $i$ to block $j$ by $A\mapsto eA$. The factor in block $j$ acts on that
strip, which is its reverse strip, by $A\mapsto Ae^{-1}$. Factors in other blocks commute with $X_{ij}(A)$. So
$\hat e\,X_{ij}(A)\,\hat e^{-1}=X_{ij}(eAe^{-1})$, and $q=\hat e_1\cdots\hat e_k$ has the required property.
\end{proof}

For a word $v(t)$ in roots with entries in the polynomial algebra $B[t]$, with $t$ central, and $T\in\grpMat_b(\F_p)$, let
$W_v(T)$ be the word obtained by substituting $T$ for $t$ in each entry and replacing each root $x_{ij}(a(t))$ by the
block root $X_{ij}(a(T))$. Substitution is a ring homomorphism $B[t]\to\grpMat_b(B)$, because numerical matrices commute
with scalar matrices from $B$. So if $v(t)$ has identity matrix image, so does $W_v(T)$. For each inner index $s$ let
$\iota_s$ be the coordinate homomorphism $x_{ij}(a)\mapsto x_{(i,s),(j,s)}(a)$. The images of $\iota_s$ for different $s$
commute, since their coordinates are disjoint.

\begin{lemma}[Triangular evaluation; {\cite[Lemma~4.2]{OAI247rf}}, {\cite[Lemmas~6.5 and~6.6]{OAI247periodic}}]
\label{lem:tor-triangular}
Let $\mathcal N$ be the subgroup generated by the roots $x_{(i,s),(j,t)}(a)$ with $s<t$, where $i=j$ is allowed. Then
$\mathcal N$ maps injectively to matrices, and is normalized by every root whose two inner indices agree. If $v(t)$ has
identity matrix image and $T\in\grpMat_b(\F_p)$ is upper triangular with every diagonal entry equal to $\lambda$, then
\[
 W_v(T)=\prod_{s=0}^{b-1}\iota_s\bigl(v(\lambda)\bigr).
\]
\end{lemma}

\begin{proof}
The proof of~\cite{OAI247rf} is characteristic-free; we recall it. Call $t-s>0$ the gap of a generator of $\mathcal N$.
Two generators either commute or are consecutive, in one of the two orders, and then their commutator is a root whose gap
is the sum of the gaps, up to the sign of its entry. They are never opposite, since both gaps are positive. Let
$\mathcal N_g$ be generated by the roots of gap at least $g$; then $\mathcal N_b=1$, each $\mathcal N_g$ is normal in
$\mathcal N$, and $\mathcal N_g/\mathcal N_{g+1}$ is generated by the commuting roots of gap $g$. Collecting
successively for $g=1,\dots,b-1$ writes every element of $\mathcal N$ as an ordered product with one factor per
coordinate pair, in increasing gap. If its matrix is the identity, the entries of gap one are the collected coefficients
of gap one, so these vanish. Once all coefficients of gap less than $g$ vanish, the entries of gap $g$ are the collected
coefficients of gap $g$, since a product of two strictly upper entries has larger gap. So all coefficients vanish, and
$\mathcal N$ is faithful. A root with equal inner indices is never opposite to a generator of $\mathcal N$, and its
commutator with one is trivial or has the same gap. So it normalizes $\mathcal N$.

For an entry $a(t)$, the matrix $a(T)$ is upper triangular with diagonal entries $a(\lambda)$. Split each block root into
the commuting product of its diagonal part and its strictly upper part. The diagonal parts normalize $\mathcal N$, so all
strictly upper parts can be collected on the right. Diagonal parts at different inner indices commute. This gives
$W_v(T)=\bigl(\prod_s\iota_s(v(\lambda))\bigr)u$ with $u\in\mathcal N$. Both $W_v(T)$ and $v(\lambda)$ have identity
matrix image, hence so does $u$, and $u=1$.
\end{proof}

\subsection{Removing the wrap of a cyclic shift}

Let $C_b$ be the numerical cyclic shift, with entries $1$ in the positions $(s,s+1)$, $s\in\Z/b$, and let $S_b$ be the
same matrix with the entry in position $(b-1,0)$ removed. So $S_b$ is strictly upper triangular.

\begin{lemma}[Removing the wrap; {\cite[Lemma~4.3]{OAI247rf}}, {\cite[Lemma~6.7]{OAI247periodic}}]\label{lem:tor-wrap}
Let $B$ be a unital $\F_p$-algebra, and let $v(t)$ be a word of $m\geq1$ roots over $B[t]$ whose entries have degree at
most $d$ and whose matrix image is the identity. Put $\delta=\max\{1,d\}$. If $b>2(m+1)\delta$, then
$W_v(C_b)=W_v(S_b)$ in $\St(B)$.
\end{lemma}

\begin{proof}
The argument of~\cite{OAI247rf} uses only supports, and we recall it to show that the signs of odd characteristic do not
affect it. The $h$th factor of $W_v(C_b)$ splits exactly as $D_hF_h$, where $D_h$ is the $h$th factor of $W_v(S_b)$ and
$F_h$ collects the entries that wrap past inner index $b-1$. A term of degree $e\leq d<b$ goes from $s$ to $s+e$, so a
wrapping term has source in $\{b-\delta,\dots,b-1\}$ and target in $\{0,\dots,\delta-1\}$. For $r\geq1$ let $V_r$ be
generated by the roots, with arbitrary outer indices, from inner indices in $H_r=\{b-r\delta,\dots,b-1\}$ to inner
indices in $L_r=\{0,\dots,r\delta-1\}$. When $b>2r\delta$ these bands are disjoint, and $V_r$ is a faithful abelian strip
by Lemma~\ref{lem:tor-strips}.

We claim that $D_h^{-1}V_rD_h\subseteq V_{r+1}$ for $r\leq m$. The constituents of $D_h^{-1}$ are roots with the same
supports as those of $D_h$ and negated entries. They all go from one outer index $i$ to a different outer index $j$,
with inner displacement between $0$ and $d$, and they commute with one another. Fix a root $y$ of $V_r$. A constituent
commuting with $y$ commutes with all conjugates of $y$ by the other constituents, and may be omitted. An interacting
constituent changes the source of $y$ only if its target is that source, and then moves it to the left by at most $d$. It
changes the target of $y$ only if its source is that target, and then moves it to the right by at most $d$. These are
the consecutive relations and their reversed forms; the sign of the new entry plays no role. No constituent is opposite
to $y$, nor to a root obtained from $y$ by one endpoint move: such a root has source minus target, in inner indices, at
least $b-(2r+1)\delta+1>\delta\geq d$. Each endpoint moves at most once. A moved source has outer index $i$, while every
constituent has target outer index $j\neq i$; likewise a moved target has outer index $j$. So all roots produced lie in
$V_{r+1}$.

The exact identity
\[
 \prod_{h=1}^m(D_hF_h)=(D_1\cdots D_m)\prod_{h=1}^m\bigl(D_{>h}^{-1}F_hD_{>h}\bigr),\qquad D_{>h}=D_{h+1}\cdots D_m,
\]
with the last product in increasing order of $h$, collects the corrections without conjugating one correction by another.
Each $D_{>h}$ has at most $m$ factors, so every correction lies in $V_{m+1}$, and $W_v(C_b)=W_v(S_b)\,u$ with
$u\in V_{m+1}$. Both evaluated words have identity matrix image, hence so does $u$. As $b>2(m+1)\delta$, the strip
$V_{m+1}$ is faithful, and $u=1$.
\end{proof}

\subsection{Graded evaluation and relative torsion}

Let $R$ be as in Theorem~\ref{thm:relative-torsion}. The map $\gamma\colon R\to R[t]$, $r_d\mapsto r_dt^d$ for
$r_d\in R_d$, is a unital ring homomorphism. For $T\in\grpMat_b(\F_p)$ and $z\in\St(R)$, let $W_z(T)$ be the image of $z$
under the composite homomorphism
\[
 \St(R)\xrightarrow{\ \gamma\ }\St(R[t])\xrightarrow{\ t\mapsto T\ }\St(\grpMat_b(R))\xrightarrow{\ \text{block roots}\ }\St(R),
\]
computed with a fixed injection of the pairs $(i,s)$ into the index set. On a representative word $z$, it is computed by
applying $\gamma$ to the entries to obtain a word $z(t)$, and then forming $W_{z(t)}(T)$. Since $W_\bullet(T)$ is a
homomorphism, $W_z(T)$ depends only on the element $z$, not on its representative. If $z\in J(R)$ then
$W_z(T)\in J(R)$, and by Lemma~\ref{lem:tor-relabel} it does not depend on the injection chosen.

\begin{lemma}[Vanishing at long cyclic shifts; {\cite[Lemma~4.4]{OAI247rf}}]\label{lem:tor-vanishing}
Let $z\in J(R)$ map to $1$ in $\St(\F_p)$, and fix a representative word of $z$ whose graded word $z(t)$ has $m\geq1$
factors with entries of degree at most $d$. Then $W_z(C_b)=1$ for every integer $b>2(m+1)\max\{1,d\}$.
\end{lemma}

\begin{proof}
By Lemmas~\ref{lem:tor-wrap} and~\ref{lem:tor-triangular}, the latter with $\lambda=0$,
$W_z(C_b)=W_{z(t)}(C_b)=W_{z(t)}(S_b)=\prod_s\iota_s(z(0))$. The word $z(0)$ consists of the degree-zero parts of the entries of $z$. It
represents the image of $z$ in $\St(\F_p)$, which is trivial; so each $\iota_s(z(0))$ is trivial.
\end{proof}

\begin{proof}[Proof of Theorem~\ref{thm:relative-torsion}]
Let $z$ be such an element, fix a representative word with $m\geq1$ factors (the empty word is trivial), and let $d$
bound the degrees of the entries of $z(t)$. Choose $a$ with $b=p^a>2(m+1)\max\{1,d\}$.

Since $C_b^b=1$ and $b$ is a power of the characteristic, $(C_b-1)^b=C_b^b-1=0$. A basis of $\F_p^b$ adapted to the
kernels of the powers of the nilpotent matrix $C_b-1$ gives $P\in\GL_b(\F_p)$ such that $P^{-1}C_bP$ is upper triangular
with all diagonal entries $1$. Put $c=\det(P)^{-1}\in\F_p^\times$ and $Q=cP$. Since $c^p=c$, also $c^b=c$, so
\[
 \det Q=c^b\det P=1,\qquad Q^{-1}C_bQ=P^{-1}C_bP=:U_b.
\]
(Over $\F_2$ every invertible matrix has determinant one, and this correction is invisible.) Apply
Lemma~\ref{lem:tor-block}(ii) to $Q^{-1}\in\SL_b(\F_p)$. It gives $q$ with $qX_{ij}(A)q^{-1}=X_{ij}(Q^{-1}AQ)$; since
$Q^{-1}a(C_b)Q=a(U_b)$ for every polynomial entry $a(t)$, we get $qW_z(C_b)q^{-1}=W_z(U_b)$. The element $W_z(C_b)$ lies
in $J(R)$, which is central (Lemma~\ref{lem:tor-center}), so $W_z(C_b)=W_z(U_b)$. By Lemma~\ref{lem:tor-triangular} with
$\lambda=1$,
\[
 W_z(U_b)=\prod_{s=0}^{b-1}\iota_s(z),
\]
since $z(1)$ represents $z$. Each $\iota_s(z)$ is represented by the word for $z$ relabelled by $i\mapsto(i,s)$, so
$\iota_s(z)=z$ by Lemma~\ref{lem:tor-relabel}. Hence $W_z(C_b)=z^b$. By Lemma~\ref{lem:tor-vanishing}, $W_z(C_b)=1$.
Therefore $z^{p^a}=1$.
\end{proof}

%======================================================================
\section{Phases, detection and the proof of Theorem~B}\label{sec:detection}

In this section we prove that finite degree truncations detect every element of $\St(R_p)$, including the central kernel,
and then assemble Theorem~B. The method is that of~\cite[Sections~2, 5, 6 and~8]{OAI247rf}. Long nonzero words carry
phases modulo integers $L$. For a representative of a central element $z$ whose entries are long words, the phases turn
evaluation at the cyclic shift $C_L$ into $z^L$. If $z$ is invisible in all truncations, the same evaluation is trivial
for large $L$, by Lemma~\ref{lem:tor-vanishing}. When $L$ is prime to $p$, this together with
Theorem~\ref{thm:relative-torsion} kills $z$.

The changes for odd $p$ are as follows. Roots and the reverse commutators of long entries have order $p$, by the
identity $H^p=[x^p,y]$. The incoming strip action of a diagonal element carries a minus sign. Coefficients in $\F_p$ are
expanded into repeated roots, and the order in which the cutoffs are chosen has to be fixed. Finally, the phase moduli
must be prime to $p$. In~\cite{OAI247rf} the moduli are the products $L_0\cdots L_d$ of the encoding lengths, which are
odd. For odd $p$ we pass to the part of $L_0\cdots L_d$ prime to $p$. Phases survive reduction to any divisor of the
modulus (Lemma~\ref{lem:det-divisor}), and these parts are unbounded (Lemma~\ref{lem:det-pprime}). So a single
presentation, the one of Theorem~A, serves every prime at once.

\subsection{The algebra and its hierarchy}

We use the hierarchy of Section~\ref{sec:prelim} with the parameters $D$ and $h_0$ fixed there once and for all, so
that $h_j=h_0D^j$ and $L_j=2\cdot2^{h_j}+1$. For a prime $p$ put
\[
 R_p=A_0(\F_p)=R_{\F_p},
\]
the unital algebra over $\F_p$ of the presentation of Theorem~A. The same table serves every prime.

\begin{proposition}[The algebra $R_p$]\label{prop:det-input}
For every prime $p$:
\begin{enumerate}[label=\textup{(\alph*)}]
\item $R_p=\bigoplus_{d\geq0}R_d$ is an infinite, finitely presented, graded unital $\F_p$-algebra with $R_0=\F_p$ and
finite-dimensional homogeneous components. Its defining relations are the elementary replacements and zero declarations
of $\mathcal S_0$, all of length at most $k$.
\item Its positive part $I=\bigoplus_{d\geq1}R_d$ is matrix-nil. In particular $R_p$ is of matrix-nil type.
\item The hierarchy $(A_j(\F_p))_{j\geq0}$ satisfies \ref{H:encoding} and \ref{H:tokens} of
Proposition~\ref{prop:hierarchy}. In~\ref{H:tokens} we write $T_j$ for a length beyond which the decomposition
$w=p_w\,t_1\cdots t_r\,s_w$ exists. Cuts are recognized from single letters: a cut sits immediately before each
occurrence of a designated letter, and is counted in a word only if that letter occurs in it~\cite[Proposition~2.2(iii)]{OAI247rf}.
\end{enumerate}
\end{proposition}

\begin{proof}
(a) Finite presentation, grading and $R_0=\F_p$ come from Lemma~\ref{lem:word-basis}. Each $R_d$ is spanned by the
finitely many words of length $d$. $R_p$ is infinite by the proof of Theorem~A. (b) is Theorem~\ref{thm:matrix-nil} with
$K=\F_p$; for $p=2$ it is~\cite[Theorem~3.7]{OAI247periodic}. (c) is Proposition~\ref{prop:hierarchy}, which holds over every
field by Remark~\ref{rem:char-free}.
\end{proof}

We use the word-class basis of Lemma~\ref{lem:word-basis}. The nonzero replacement classes form a basis of $R_p$. A word
$w$ has $[w]=0$ if and only if a finite sequence of replacements turns it into a word containing a zero-declared subword.
Hence a subword of a nonzero word is nonzero, and the product of two basis classes is either zero or the class of their
concatenation. We use the same symbol for a nonzero word and its basis class when it occurs as an entry. Each
$R_{\geq d}=\bigoplus_{e\geq d}R_e$ is a two-sided ideal, and $R_p/R_{\geq d}$ is finite.

\subsection{Phases}

\begin{definition}\label{def:phase-system}
Let $L,N\geq1$. A \emph{phase system modulo $L$ with threshold $N$} assigns to every nonzero word $w$ over $\Sigma_0$
with $|w|\geq N$ a \emph{start phase} $\sigma(w)\in\Z/L$, with \emph{end phase} $\tau(w)=\sigma(w)+|w|$, such that:
\begin{enumerate}[label=\textup{(P\arabic*)}]
\item $\sigma(w)$ depends only on the class $[w]$;
\item if $v$ is a subword of $w$ with $|v|\geq N$, starting at offset $b$ from the beginning of $w$, then
$\sigma(v)=\sigma(w)+b$.
\end{enumerate}
\end{definition}

\begin{lemma}[Phases of long words; {\cite[Lemma~2.3]{OAI247rf}}]\label{lem:det-phases}
For every $d\geq0$ there is a phase system modulo $L_0L_1\cdots L_d$ with some threshold $N$.
\end{lemma}

The proof in~\cite{OAI247rf} uses only the hierarchy and the word-class basis. Both are characteristic-free, so the
proof applies verbatim over $\F_p$. We give it, since it is the source of all the phase information.

\begin{proof}
Put $Q_j=L_jL_{j+1}\cdots L_d$ for $j\leq d$ and $Q_{d+1}=1$. We construct, by backward induction on $j=d+1,d,\dots,0$,
\emph{literal} phases $\sigma_j(w)\in\Z/Q_j$ for nonzero words $w$ over $\Sigma_j$ with $|w|\geq n_j$, satisfying (P2) but
not yet (P1). At level $d+1$ take $\sigma_{d+1}=0$ and $n_{d+1}=1$.

Let $j\leq d$. For each token $t\in\mathcal T_j$ fix one word $u_t$ as in~\ref{H:tokens}, the same at every
occurrence. For a nonzero word $w$ with its decomposition
$w=p_wt_1\cdots t_rs_w$, put $U(w)=u_{t_1}\cdots u_{t_r}$, and choose
\[
 n_j\geq\max\{1,\,T_j,\,2B_j+L_jn_{j+1}\}.
\]
If $|w|\geq n_j$, then $|U(w)|=(|w|-|p_w|-|s_w|)/L_j\geq n_{j+1}$. Moreover $U(w)$ is nonzero: the homomorphism $E_j$
and the token equalities give $[w]=[p_w]\,E_j([U(w)])\,[s_w]$, so $[U(w)]=0$ would force $[w]=0$. No injectivity of $E_j$
is used. Define
\[
 \sigma_j(w)=L_j\,\sigma_{j+1}(U(w))-|p_w|\pmod{Q_j}.
\]
This is well defined because multiplication by $L_j$ carries $\Z/Q_{j+1}$ to $\Z/L_jQ_{j+1}=\Z/Q_j$.

Let $v$ be a subword of $w$ with $|v|\geq n_j$, starting at offset $b$. Cuts are recognized from single letters, so the
cuts of $v$ are exactly the cuts of $w$ whose following letters lie in $v$. Hence the complete tokens of $v$ are a
consecutive run $t_\alpha,\dots,t_\beta$ of the tokens of $w$, and with the fixed choices, $U(v)$ is a literal subword of
$U(w)$, at upper offset $a=|u_{t_1}|+\dots+|u_{t_{\alpha-1}}|$. The first complete token of $v$ starts in $w$ at position
\[
 b+|p_v|=|p_w|+L_ja.
\]
Applying (P2) at level $j+1$ and then this equality,
\[
 \sigma_j(v)=L_j\bigl(\sigma_{j+1}(U(w))+a\bigr)-|p_v|=L_j\sigma_{j+1}(U(w))-|p_w|+b=\sigma_j(w)+b\pmod{Q_j}.
\]
This completes the induction. It has finitely many steps, and no phase at infinite depth is needed.

At level $0$ we have literal phases modulo $Q_0=L_0\cdots L_d$ with (P2) for lengths at least $n_0$. Put $N=2n_0+k$.
A single replacement inside a nonzero word of length at least $N$ changes at most $k$ consecutive letters and preserves
length. So an unchanged prefix or an unchanged suffix of length at least $n_0$ is common to both words, at the same
offset. Applying (P2) to it in both words shows that their phases agree. Every word in a nonzero class is nonzero, and
any two words in the class are joined by a finite sequence of replacements; so the phase is constant on the class. This
is (P1). Property (P2) persists when the threshold is raised from $n_0$ to $N$.
\end{proof}

\begin{lemma}[Reduction to a divisor]\label{lem:det-divisor}
Let $\sigma$ be a phase system modulo $L$ with threshold $N$, and let $L'$ divide $L$. Then reducing $\sigma$ modulo $L'$
gives a phase system modulo $L'$ with threshold $N$.
\end{lemma}

\begin{proof}
Properties (P1) and (P2) are congruences modulo $L$, and they imply the same congruences modulo $L'$.
\end{proof}

\begin{corollary}[{\cite[Corollary~2.4]{OAI247rf}}]\label{cor:det-products}
Fix a phase system modulo $L$ with threshold $N$. If $u,v$ are nonzero words of lengths at least $N$ and $uv\neq0$,
then $\tau(u)=\sigma(v)$, $\sigma(uv)=\sigma(u)$ and $\tau(uv)=\tau(v)$. In particular, $\tau(u)\neq\sigma(v)$ implies
$uv=0$.
\end{corollary}

\begin{proof}
Apply (P2) to the prefix $u$ and the suffix $v$ of the nonzero word $uv$: $\sigma(u)=\sigma(uv)$ and
$\sigma(v)=\sigma(uv)+|u|=\tau(u)$. Adding $|v|$ gives the last identity.
\end{proof}

For a positive integer $x$, write $x_{(p')}$ for its largest divisor prime to $p$.

\begin{lemma}[The parts prime to $p$ are unbounded]\label{lem:det-pprime}
For every prime $p$, $(L_0L_1\cdots L_d)_{(p')}\to\infty$ as $d\to\infty$.
\end{lemma}

\begin{proof}
Since $(L_0\cdots L_d)_{(p')}\geq(L_d)_{(p')}$, it suffices to show $(L_d)_{(p')}\to\infty$. For $p=2$, each $L_j$ is odd
and $L_j\to\infty$. Let $p$ be odd, let $o$ be the order of $2$ modulo $p$, and put $C=v_p(2^o-1)\geq1$. We first show
\begin{equation}\label{eq:det-lte}
 v_p(2^n+1)\leq C+v_p(n)\qquad(n\geq1).
\end{equation}
If $p\nmid2^n+1$ there is nothing to prove. Otherwise $2^{2n}\equiv1$, so $o$ divides $2n$, and
$v_p(2^n+1)\leq v_p(2^{2n}-1)=v_p(a^e-1)$ with $a=2^o$ and $e=2n/o$. We use the elementary half of the lifting-the-exponent
lemma: if $a\equiv1\pmod p$ then $v_p(a^e-1)=v_p(a-1)+v_p(e)$. Indeed, if $p\nmid e'$ and $x\equiv1\pmod p$, then
$(x^{e'}-1)/(x-1)=\sum_{i<e'}x^i\equiv e'\not\equiv0$, so $v_p(x^{e'}-1)=v_p(x-1)$. And if $x=1+y$ with $v_p(y)=s\geq1$,
then in $x^p-1=\sum_{i=1}^p\binom pi y^i$ the term $py$ has valuation $s+1$, the terms with $2\leq i\leq p-1$ have
valuation at least $1+2s$, and $y^p$ has valuation $ps\geq s+2$, because $p\geq3$. So $v_p(x^p-1)=s+1$. Writing
$e=p^fe'$ with $p\nmid e'$ and iterating gives the formula. Since $p$ is odd, $v_p(e)\leq v_p(2n)=v_p(n)$,
proving~\eqref{eq:det-lte}.

Now $L_d=2^{h_d+1}+1$, so by~\eqref{eq:det-lte}, $p^{v_p(L_d)}\leq p^C(h_d+1)$ and
\[
 (L_d)_{(p')}\geq\frac{2^{h_d+1}+1}{p^C\,(h_d+1)},
\]
which tends to infinity with $h_d$.
\end{proof}

\begin{remark}\label{rem:det-coprime}
Instead of passing to divisors, one can choose the parameters so that every $L_j$ is prime to $p$, using the freedom of
Remark~\ref{rem:parameters}. For $p\geq5$, take $h_0$ divisible by the order of $2$ modulo $p$; then
$2^{h_j}\equiv1$ and $L_j\equiv3\pmod p$. For $p=3$, take $D$ and $h_0$ odd; then each $h_j$ is odd and
$L_j\equiv-1\pmod3$. This route makes the table depend on $p$. The divisor reduction above avoids that, and is the route
we follow.
\end{remark}

\subsection{Opposite roots and long entries}

From now until Lemma~\ref{lem:det-high}, fix a phase system modulo $L$ with threshold $N$, let $k$ bound the lengths of
the defining relations, and fix an integer
\[
 M\geq2N+k+1.
\]
All groups are stable Steinberg groups over $R=R_p$. Lengths, phases and products of nonzero basis classes are as above.

\begin{lemma}[Opposite roots; after {\cite[Lemma~5.1]{OAI247rf}}]\label{lem:det-opposite}
Let $a,c\in R$ with $ca=0$. For distinct indices $i,h,j$ and every $b\in R$,
\begin{equation}\label{eq:det-cut}
 [x_{ij}(ab),x_{ji}(c)]=[x_{ih}(a),x_{hi}(bc)].
\end{equation}
In particular $H=[x_{ij}(a),x_{ji}(c)]$ does not depend on $j\neq i$. It commutes with $x_{ij}(a)$ and $x_{ji}(c)$, it
satisfies $H^p=1$, and its matrix is $1+acE_{ii}$.
\end{lemma}

\begin{proof}
Apply Lemma~\ref{lem:three-factor} to $X=x_{ih}(a)$, $Y=x_{hj}(b)$ and $Z=x_{ji}(c)$. Then $A=[X,Y]=x_{ij}(ab)$ and
$B=[Y,Z]=x_{hi}(bc)$. The pairs $(A,X)$ and $(B,Z)$ are in commuting positions. Also $[X,Z]=1$, because
$[x_{ji}(c),x_{ih}(a)]=x_{jh}(ca)=1$, and $[A,B]=1$, because $[x_{hi}(bc),x_{ij}(ab)]=x_{hj}(bcab)=x_{hj}(b(ca)b)=1$. The
lemma gives~\eqref{eq:det-cut}; with $b=1$ it gives independence of the auxiliary index.

Write $H=[x_{ih}(a),x_{hi}(c)]$ with $h\notin\{i,j\}$. The factor $x_{ih}(a)$ commutes with $x_{ij}(a)$, being in a
commuting position, and with $x_{ji}(c)$, because $[x_{ji}(c),x_{ih}(a)]=x_{jh}(ca)=1$. Likewise $x_{hi}(c)$ commutes with
$x_{ji}(c)$, and with $x_{ij}(a)$ because $[x_{hi}(c),x_{ij}(a)]=x_{hj}(ca)=1$. So $H$ commutes with $X_0=x_{ij}(a)$ and
$Z_0=x_{ji}(c)$. As in Step~3 of the proof of Theorem~\ref{thm:fp}, $[X_0^p,Z_0]=H^p$, and $X_0^p=x_{ij}(pa)=1$. So
$H^p=1$. In characteristic two~\cite{OAI247rf} uses instead that each root is an involution. Multiplying the four
elementary matrices, using $ca=0$, gives $1+acE_{ii}$.
\end{proof}

\begin{lemma}[Long opposite commutators; {\cite[Lemma~5.2]{OAI247rf}}]\label{lem:det-long}
For each index $i$ and each nonzero class $w$ with $|w|\geq M$ and $\sigma(w)\neq\tau(w)$, there is $H_i(w)\in\St(R)$
with the following property. If $u,q$ are nonzero words with
\[
 |u|,|q|\geq N,\qquad|u|+|q|\geq M,\qquad\sigma(u)\neq\tau(q),
\]
then for every $j\neq i$,
\[
 [x_{ij}(u),x_{ji}(q)]=\begin{cases}H_i(w)&\text{if }uq\text{ represents the nonzero class }w,\\ 1&\text{if }uq=0.\end{cases}
\]
Moreover $H_i(w)^p=1$, and its matrix is $1+wE_{ii}$.
\end{lemma}

\begin{proof}
This is the argument of~\cite{OAI247rf}, with Lemma~\ref{lem:det-opposite} in place of its characteristic-two version.
Since $\tau(q)\neq\sigma(u)$, Corollary~\ref{cor:det-products} gives $qu=0$, so Lemma~\ref{lem:det-opposite} applies.

\emph{Moving the cut.} Consider the literal word $uq$ cut after $u$. To move the cut to the right, write $q=q'c$ with
$|c|\geq N$. Then $c$ is nonzero, $\tau(c)=\tau(q)\neq\sigma(u)$, and so $cu=0$. Now~\eqref{eq:det-cut}, with independence
of the auxiliary index, gives $[x_{ij}(uq'),x_{ji}(c)]=[x_{ij}(u),x_{ji}(q)]$. If $uq'=0$ this common value is $1$.
Otherwise $\sigma(uq')=\sigma(u)$, and the new cut keeps the mismatch. Moves to the left are symmetric: if $u=au'$ with
$|a|\geq N$, then $\sigma(a)=\sigma(u)$, $qa=0$, and $[x_{ij}(au'),x_{ji}(q)]=[x_{ij}(a),x_{ji}(u'q)]$. So the commutator is
unchanged by any move of the cut within the positions $[N,|uq|-N]$, and if a part becomes zero, the commutator is $1$.

\emph{Replacements.} Apply one replacement, of length $\ell\leq k$, inside $uq$. The allowed cut positions $[N,|uq|-N]$
number at least $M-2N+1\geq k+2$, while the replacement has at most $\ell-1\leq k-1$ of them strictly inside it. Move the
cut to an allowed position outside the replacement, which then lies in one part. That part keeps its class, and if
nonzero, its phases; the other part is unchanged. So the commutator and the mismatch survive the replacement.

If $uq$ represents a nonzero class $w$, all its subwords and all words in its class are nonzero, and the two steps show
that the commutator depends only on $w$. Every such $w$ can be cut into parts of lengths at least $N$, and these have
start phase $\sigma(w)$ and end phase $\tau(w)$ by (P2). Define $H_i(w)$ by any such cut. If $uq=0$, follow a sequence of
replacements to a word containing a zero-declared subword, which has length at most $k$. Either a part has already
become zero, or the cut can be moved so that the zero-declared subword lies in one part. That part is zero, and the
commutator is $1$. The order and the matrix of $H_i(w)$ come from Lemma~\ref{lem:det-opposite}.
\end{proof}

Let $S$ be a finite set of indices with a designated phase $\phi(i)\in\Z/L$ for each $i\in S$. An entry $w$ from $i$ to
$j$ is \emph{consistent} if $w$ is a nonzero basis class with $|w|\geq M$, $\sigma(w)=\phi(i)$ and $\tau(w)=\phi(j)$. It
is a \emph{wrong-row entry} if instead $|w|\geq M$, $\sigma(w)\neq\phi(i)$ and $\tau(w)=\phi(j)$. Let $\mathcal U(S,\phi)$
be the subgroup generated by the roots $x_{ij}(w)$, $i\neq j$, and the elements $H_i(w)$, $i=j$, for all wrong-row
entries $w$.

\begin{lemma}[The wrong-row subgroup; after {\cite[Lemma~5.3]{OAI247rf}}]\label{lem:det-wrong-row}
The subgroup $\mathcal U(S,\phi)$ is abelian of exponent dividing $p$, and the matrix map is injective on it. Every root
$x_{ef}(b)$ with a consistent entry $b$ normalizes it.
\end{lemma}

\begin{proof}
Write $\mathcal U=\mathcal U(S,\phi)$. Let $b$ be a wrong-row or consistent entry from $e$ to $f$. In either case
$\tau(b)=\phi(f)$.

\emph{Off-diagonal generators.} Let $x_{ij}(a)$ be an off-diagonal generator. If $j=e$ and $i\neq f$, then
$x_{jf}(b)\,x_{ij}(a)\,x_{jf}(b)^{-1}=x_{if}(-ab)\,x_{ij}(a)$. If $b$ is wrong-row, then $ab=0$, since
$\tau(a)=\phi(j)\neq\sigma(b)$. If $b$ is consistent and $ab\neq0$, then $ab$ has $\sigma(ab)=\sigma(a)\neq\phi(i)$ and
$\tau(ab)=\tau(b)=\phi(f)$, so it is a wrong-row entry from $i$ to $f$, and $x_{if}(-ab)=x_{if}(ab)^{p-1}\in\mathcal U$.
If $f=i$ and $j\neq e$, then $ba=0$, since $\tau(b)=\phi(i)\neq\sigma(a)$, and the roots commute. If the positions are
opposite, $e=j$ and $f=i$, then $\sigma(a)\neq\phi(i)=\tau(b)$, and Lemma~\ref{lem:det-long} gives
$[x_{ij}(a),x_{ji}(b)]=H_i(ab)$ if $ab\neq0$, and $1$ otherwise. In the first case $ab$ is a wrong-row entry from $i$ to
$i$, since $\tau(ab)=\tau(b)=\phi(i)$. If $b$ is wrong-row, then $\sigma(b)\neq\phi(j)=\tau(a)$, so $ab=0$ and the roots
commute. All other pairs of positions are commuting positions.

\emph{Diagonal generators.} Represent $H_i(w)$ as a commutator $[x_{ih}(a),x_{hi}(c)]$ with $ac=w$, $|a|,|c|\geq N$ and an
auxiliary index $h\notin S\cup\{e,f\}$. If $i\notin\{e,f\}$, the supports are disjoint and the elements commute. Otherwise
$x_{ef}(b)$ lies in a strip between $\{i,h\}$ and a disjoint set. The matrix of $H_i(w)$ is $1+wE_{ii}$, with inverse
$1-wE_{ii}$ since $w^2=0$. So Lemma~\ref{lem:tor-strips} gives, for $f\neq i$ and $e\neq i$,
\[
 H_i(w)\,x_{if}(b)\,H_i(w)^{-1}=x_{if}(b+wb),\qquad H_i(w)\,x_{ei}(b)\,H_i(w)^{-1}=x_{ei}(b-bw).
\]
Over $\F_2$ the two signs agree, and~\cite{OAI247rf} writes a plus sign in both. Over $\F_p$ the sign in the incoming
action matters only formally, because $bw=0$ always, as $\tau(b)=\phi(i)\neq\sigma(w)$. If $b$ is wrong-row then also $wb=0$, as $\tau(w)=\phi(i)\neq\sigma(b)$.
If $b$ is consistent and $wb\neq0$, then $wb$ is a wrong-row entry from $i$ to $f$, and $H_i(w)\,x_{if}(b)\,H_i(w)^{-1}=
x_{if}(wb)\,x_{if}(b)$; hence $x_{if}(b)\,H_i(w)\,x_{if}(b)^{-1}=x_{if}(wb)^{-1}H_i(w)\in\mathcal U$. Two diagonal
generators at different indices are represented with disjoint auxiliary pairs, and commute. For two at the same index,
write $H_i(w')=[x_{it}(a),x_{ti}(c)]$ with $ac=w'$ and $|a|,|c|\geq N$, using an index $t$ different from the auxiliary
index of $H_i(w)$. Then $\sigma(a)=\sigma(w')\neq\phi(i)=\tau(w)$ and $\tau(c)=\tau(w')=\phi(i)\neq\sigma(w)$, so $wa=0$ and
$cw=0$. The displayed strip actions then show that $H_i(w)$ fixes $x_{it}(a)$ and $x_{ti}(c)$, so it commutes with
$H_i(w')$.

\emph{Conclusion.} All generators of $\mathcal U$ commute. Each has order dividing $p$: $x_{ij}(w)^p=x_{ij}(pw)=1$, and
$H_i(w)^p=1$ by Lemma~\ref{lem:det-long}. So every element of $\mathcal U$ is a product of distinct generators with
exponents in $\{0,\dots,p-1\}$, and its matrix is
\[
 1+\sum_{(i,j,w)}c_{ij,w}\,w\,E_{ij}.
\]
All cross products of the added matrices vanish: a nonzero product would need the end phase $\phi(j)$ of one entry to
equal the start phase of the next, which differs from the designated phase $\phi(j)$ of its row. As the nonzero classes
are linearly independent, the matrix is the identity only if all exponents vanish. So $\mathcal U$ maps injectively to
matrices. Conjugation by a consistent root $x_{ef}(b)$ sends every generator into $\mathcal U$. Since
$x_{ef}(b)^{-1}=x_{ef}(b)^{p-1}$, conjugation by the inverse does so too, and $x_{ef}(b)$ normalizes $\mathcal U$. In
characteristic two, \cite{OAI247rf} uses here that each root is an involution.
\end{proof}

\subsection{Relocation and the \texorpdfstring{$L$}{L}th power}

\begin{lemma}[Relocation by phases; {\cite[Lemma~6.1]{OAI247rf}}]\label{lem:det-relocation}
Let $z\in J(R)$ be represented by a word
\[
 z=\prod_{\nu=1}^m x_{i_\nu j_\nu}(w_\nu)
\]
whose entries $w_\nu$ are nonzero basis classes with $|w_\nu|\geq M$. Label new coordinates by $(i,s)$, with $s\in\Z/L$.
Then
\begin{equation}\label{eq:det-relocation}
 z=\prod_{\nu=1}^m x_{(i_\nu,\sigma(w_\nu)),(j_\nu,\tau(w_\nu))}(w_\nu)\qquad\text{in }\St(R).
\end{equation}
\end{lemma}

The hypothesis asks for basis classes with coefficient one. A root $x_{ij}(cw)$ with $c\in\F_p$ is the power
$x_{ij}(w)^c$, $0\leq c<p$, and an inverse root $x_{ij}(w)^{-1}$ is $x_{ij}(w)^{p-1}$. So any word whose entries are
multiples of long basis classes can be rewritten in this form. This changes the number of factors, but not the degrees.

\begin{proof}
Assume $m\geq1$, and let $S$ be a finite set of at least three indices containing the original ones. Put $A=\F_p\oplus R_{\geq M}$, a
subalgebra of $R$. Choose numerical columns $u_s\in\F_p^L$ and rows $v_s\in(\F_p^L)^*$ with $v_su_s=1$ for all
$s\in\Z/L$, and define $\rho\colon A\to\grpMat_L(R)$ on the basis by
\[
 \rho(w)=u_{\sigma(w)}v_{\tau(w)}\,w,\qquad\rho(1)=1,
\]
extended $\F_p$-linearly. This is a unital algebra homomorphism. If $ww'\neq0$, then $\tau(w)=\sigma(w')$ and $ww'$ has
the outer phases $\sigma(w)$ and $\tau(w')$, by Corollary~\ref{cor:det-products}, so
$\rho(w)\rho(w')=u_{\sigma(w)}(v_{\tau(w)}u_{\sigma(w')})v_{\tau(w')}ww'=\rho(ww')$. If $ww'=0$, both sides vanish
whatever the middle pairing is. The matrix of the given word is the identity over $A$, since $A\subseteq R$. Applying
$\rho$ to the entries and passing to block roots therefore gives a word with identity matrix image, which represents an
element of $J(R)$. We change $\rho$ in three steps; the evaluated element stays equal to $z$ throughout.

\emph{Start.} Take $u_s=\grpbe_0$ and $v_s=\grpbe_0^*$ for all $s$. The evaluated word is the given word relabelled by
$i\mapsto(i,0)$, so it represents $z$ by Lemma~\ref{lem:tor-relabel}. The characteristic-two lemma used here
in~\cite{OAI247rf} is replaced by that lemma.

\emph{Unused columns.} Keep $u_s=\grpbe_0$, and change the rows to $v_0=\grpbe_0^*$ and $v_t=\grpbe_0^*+\grpbe_t^*$ for
$t\neq0$. Every pairing $v_tu_s$ is still $1$. Each block root acquires at most one extra factor, from a coordinate
$(i,0)$ to an unused coordinate $(j,t)$ with $t\neq0$. These factors lie in the strip from $S\times\{0\}$ to
$S\times(\Z/L\setminus\{0\})$, which is normalized by the original factors, supported on $S\times\{0\}$
(Lemma~\ref{lem:tor-strips}). Collecting them writes the new word as the old one times an element of this strip. That
element has identity matrix image, so it is trivial.

\emph{Change of basis.} Let $Q$ be the numerical matrix with rows $v_t$. Then $Q=\prod_{t\neq0}e_{t0}(1)$ is unipotent
and elementary, of determinant one, with $Q\grpbe_0=\sum_r\grpbe_r$ and $v_tQ^{-1}=\grpbe_t^*$. Conjugating every entry
matrix by $Q$ changes the columns to $u_s=\sum_r\grpbe_r$ and the rows to $v_t=\grpbe_t^*$. By
Lemma~\ref{lem:tor-block}(ii), this simultaneous change is induced by conjugation in $\St(R)$. The evaluated element is
central, so it is unchanged.

\emph{Wrong rows.} Give $(i,s)$ the designated phase $\phi(i,s)=s$. A block root with entry $w$ is now
$\prod_{r\in\Z/L}x_{(i,r),(j,\tau(w))}(w)$, with commuting factors. The factor with $r=\sigma(w)$ is consistent and is the
factor kept in~\eqref{eq:det-relocation}; the others are wrong-row generators of $\mathcal U(S\times\Z/L,\phi)$.
Consistent roots normalize this subgroup (Lemma~\ref{lem:det-wrong-row}), so collecting the wrong-row factors writes the
evaluated word as the right side of~\eqref{eq:det-relocation} times some $u\in\mathcal U(S\times\Z/L,\phi)$. The right
side of~\eqref{eq:det-relocation} has identity matrix image: it is the evaluation by $\rho$ with $u_s=\grpbe_s$ and
$v_t=\grpbe_t^*$, which is again a homomorphism. So $u$ has identity matrix image, and $u=1$ by
Lemma~\ref{lem:det-wrong-row}.
\end{proof}

\begin{proposition}[{\cite[Proposition~6.2]{OAI247rf}}]\label{prop:det-power}
Under the hypotheses of Lemma~\ref{lem:det-relocation}, $W_z(C_L)=z^L$.
\end{proposition}

\begin{proof}
For a basis class $w$ of degree $|w|$, the cyclic evaluation of $x_{ij}(w)$ is the block root with entry $wC_L^{|w|}$.
Since $\tau(w)-\sigma(w)=|w|$ in $\Z/L$, its factors are
\[
 x_{(i,\sigma(w)+a),(j,\tau(w)+a)}(w),\qquad a\in\Z/L.
\]
Factors with different offsets $a\neq a'$ commute, even when they come from different roots. Let their entries be $w$
and $w'$. If the target of the first is the source of the second, then $\tau(w)+a=\sigma(w')+a'$, so
$\tau(w)\neq\sigma(w')$ and $ww'=0$. Symmetrically, the reverse coincidence forces $w'w=0$. So the Steinberg relations
give commutation unless the two roots are opposite. In that case both products vanish, and $\sigma(w)\neq\tau(w')$, so
Lemma~\ref{lem:det-long} gives commutation. We may therefore collect $W_z(C_L)$, computed on the given representative, by
offsets, keeping the original order within each offset. The offset-zero word is the right side
of~\eqref{eq:det-relocation}, which is $z$ by Lemma~\ref{lem:det-relocation}. The offset-$a$ word is its relabelling by
$(i,s)\mapsto(i,s+a)$, so it also equals $z$, by Lemma~\ref{lem:tor-relabel}. In characteristic two~\cite{OAI247rf} uses
the permutation lemma here. There are $L$ offsets.
\end{proof}

\subsection{Detection by degree truncations}

\begin{lemma}[High entries; after {\cite[Lemma~6.3]{OAI247rf}}]\label{lem:det-high}
For every $M\geq1$, each element of $\ker\bigl(\St(R)\to\St(R/R_{\geq2M})\bigr)$ is represented by a finite product of
roots $x_{ij}(w)$ whose entries $w$ are nonzero basis classes with $|w|\geq M$.
\end{lemma}

\begin{proof}
For a two-sided ideal $K$, the kernel of $\St(R)\to\St(R/K)$ is the normal closure of the roots $x_{ij}(a)$, $a\in K$:
after these are killed, $x_{ij}(a)$ depends only on $a$ modulo $K$, and the remaining relations are those of
$\St(R/K)$. For $K=R_{\geq2M}$, expand each entry in the basis with coefficients in $\{0,\dots,p-1\}$, which turns a root
into a product of positive powers of roots $x_{ij}(w)$ with $|w|\geq2M$. Consider a conjugate $gx_{ij}(w)g^{-1}$. Write
$w=ab$ as words with $|a|,|b|\geq M$; both are nonzero. Let $S$ contain $i,j$ and the indices of a word for $g$, and
choose $h\notin S$. Then
\[
 gx_{ij}(w)g^{-1}=[gx_{ih}(a)g^{-1},\,gx_{hj}(b)g^{-1}],
\]
and by Lemma~\ref{lem:tor-strips}, with $A_g$ the matrix of $g$ on $S$,
\[
 gx_{ih}(a)g^{-1}=\prod_{r\in S}x_{rh}\bigl((A_g)_{ri}\,a\bigr),\qquad
 gx_{hj}(b)g^{-1}=\prod_{s\in S}x_{hs}\bigl(b\,(A_g^{-1})_{js}\bigr).
\]
All these entries lie in the ideal $R_{\geq M}$. Expand them in the basis as above, and expand the commutator as a
product, writing each inverse root $x(w')^{-1}$ as $x(w')^{p-1}$. Every entry is then a basis class of length at least
$M$.
\end{proof}

\begin{theorem}[Detection; {\cite[Theorem~6.4]{OAI247rf}} over $\F_p$]\label{thm:det-detection}
Degree truncations detect the stable Steinberg group of $R=R_p$:
\[
 \bigcap_{d\geq1}\ker\bigl(\St(R)\longrightarrow\St(R/R_{\geq d})\bigr)=1.
\]
\end{theorem}

\begin{proof}
Let $z$ lie in the intersection. Its matrix is the identity modulo every $R_{\geq d}$, and $\bigcap_dR_{\geq d}=0$, so
$z\in J(R)$. Taking $d=1$, $z$ maps to $1$ in $\St(\F_p)$.

The order of choices matters. First fix one representative word of $z$, the \emph{original} representative; if it is
empty, $z=1$. Let $m\geq1$ be the number of factors of its graded word and $d$ the largest degree of its entries. By Lemma~\ref{lem:tor-vanishing},
$W_z(C_b)=1$ for every $b>b_0=2(m+1)\max\{1,d\}$. By Lemma~\ref{lem:det-pprime} there is $e$ with
\[
 L=(L_0L_1\cdots L_e)_{(p')}>b_0,
\]
and Lemmas~\ref{lem:det-phases} and~\ref{lem:det-divisor} give a phase system modulo $L$ with some threshold $N$. Only
now choose $M\geq2N+k+1$. Since $z$ dies in $\St(R/R_{\geq2M})$, Lemma~\ref{lem:det-high}, followed by the expansion of
coefficients described after Lemma~\ref{lem:det-relocation}, gives a second representative whose entries are basis
classes of length at least $M$. Its number of factors may be very large. This does not affect the choice of $L$:
$W_z(C_L)$ is the value at $z$ of a homomorphism, so it does not depend on the representative, and
Proposition~\ref{prop:det-power} imposes no bound relating $L$ to the second representative. Applying
Proposition~\ref{prop:det-power} to it,
\[
 z^L=W_z(C_L)=1.
\]
By Theorem~\ref{thm:relative-torsion}, $z^{p^a}=1$ for some $a$. As $L$ is prime to $p$, $z=1$.
\end{proof}

\subsection{Finite quotients and the proof of Theorem~B}

\begin{proposition}[Finite truncations; {\cite[Proposition~8.1]{OAI247rf}} over $\F_p$]\label{prop:det-truncations}
For every $d\geq1$ the group $\St_{12}(R_p/R_{\geq d})$ is finite, and the homomorphisms
$\St_{12}(R_p)\to\St_{12}(R_p/R_{\geq d})$ separate the elements of $\St_{12}(R_p)$.
\end{proposition}

\begin{proof}
Put $B_d=R_p/R_{\geq d}$. It is finite, being the sum of finitely many finite-dimensional homogeneous components. Its
positive part $I_d$ satisfies $I_d^d=0$, so every square matrix over $I_d$ is nilpotent, and $B_d$ is of matrix-nil
type. By Proposition~\ref{prop:centrality}, the kernel of $\St_{12}(B_d)\to\E_{12}(B_d)$ is central. The matrix group
$\E_{12}(B_d)$ is finite, so the center of $\St_{12}(B_d)$ has finite index. By Schur's theorem~\cite{Schur1904},
\cite[Theorem~10.1.4]{Robinson1996}, the derived subgroup is finite. As $\St_{12}(B_d)$ is perfect
(Lemma~\ref{lem:fp-perfect}), it is finite.

Let $g\neq1$ in $\St_{12}(R_p)$. Its image in $\St(R_p)$ is nontrivial by Theorem~\ref{thm:stability}, and it is detected
in $\St(B_d)$ for some $d$ by Theorem~\ref{thm:det-detection}. The square formed by $\St_{12}(R_p)\to\St(R_p)$,
$\St_{12}(B_d)\to\St(B_d)$ and the two truncation maps commutes on root generators. So the image of $g$ in
$\St_{12}(B_d)$ is nontrivial. No injectivity of $\St_{12}(B_d)\to\St(B_d)$ is needed.
\end{proof}

\begin{proof}[Proof of Theorem~B]
Let $R_p$ be as in Proposition~\ref{prop:det-input}; it is $A_0(\F_p)$ for the presentation of Theorem~A. It is a
finitely presented graded $\F_p$-algebra $\F_p\oplus I$ with finite-dimensional homogeneous components and matrix-nil
positive part.

(i) Put $\Gamma_p=\St_{12}(R_p)$. It has an ordinary finite presentation by Theorem~\ref{thm:fp}, with $n=12$. It is
infinite, because $R_p$ is infinite (Lemma~\ref{lem:fp-perfect}). It is residually finite by
Proposition~\ref{prop:det-truncations}. Periodicity follows from (ii): the image of $\gamma\in\Gamma_p$ in the finite
group $\St_{12}(\F_p)$ has some finite order $e$, so $\gamma^e\in G_p$, and a $p$-power of $\gamma^e$ is trivial.

(ii) Put $G_p=\ker(\Gamma_p\to\St_{12}(\F_p))$. The target is finite by Proposition~\ref{prop:det-truncations} with
$d=1$, so $G_p$ has finite index in $\Gamma_p$. Hence $G_p$ is infinite, it has an ordinary finite presentation by the
Reidemeister--Schreier theorem~\cite[Chapter~II, Section~4]{LyndonSchupp1977}, and it is residually finite. Let
$g\in G_p$. Its matrix reduces to the identity over $\F_p$, so it equals $1+A$ with $A\in M_{12}(I)$. Lemma~\ref{lem:fp-local}(iv)
gives $a$ with $(1+A)^{p^a}=1$, so $g^{p^a}\in J_{12}(R_p)$, and $g^{p^a}$ still maps to $1$ in $\St_{12}(\F_p)$, hence in
$\St(\F_p)$. Its image in $\St(R_p)$ therefore lies in $J(R_p)\cap\ker(\St(R_p)\to\St(\F_p))$. By
Theorem~\ref{thm:relative-torsion} that image is killed by some $p^b$. By Theorem~\ref{thm:stability},
$\Gamma_p\to\St(R_p)$ is injective, so $g^{p^{a+b}}=1$. Thus every element of $G_p$ has finite $p$-power order.
\end{proof}

No statement about the torsion of $J_{12}(\F_p)$, and no inference from residual finiteness of $\E_{12}(R_p)$ to that of
$\Gamma_p$, has been used. The group $\Gamma_p$ is periodic; the theorem does not assert that it is a $p$-group.

\begin{corollary}[Property (T)]\label{cor:det-T}
The groups $\Gamma_p$ and $G_p$ have Kazhdan's property~(T).
\end{corollary}

\begin{proof}
Since $\F_p$ is a prime field, $R_p$ is generated as a ring by $1$ and its finitely many letters. By Ershov and
Jaikin-Zapirain~\cite[Theorem~6.2]{ErshovJaikinZapirain2010}, $\St_n(S)$ has property~(T) for every finitely generated
unital ring $S$ and every $n\geq3$; so $\Gamma_p$ has property~(T). Property~(T) passes to subgroups of finite
index~\cite[Section~1.7]{BekkaHarpeValette2008}, so $G_p$ has it too. As in~\cite[Section~6.5]{OAI247periodic}, it follows
that $\Gamma_p$ and $G_p$ are not amenable.
\end{proof}

\begin{corollary}[Unbounded exponent]\label{cor:det-exponent}
The orders of the elements of $G_p$ are unbounded.
\end{corollary}

\begin{proof}
Suppose every element of $G_p$ had order dividing $p^a$. The group $G_p$ is finitely generated, say by $r$ elements, so
each of its finite quotients is an $r$-generator finite group of exponent dividing $p^a$. By Zel'manov's solution of the
restricted Burnside problem~\cite{ZelmanovOdd,ZelmanovTwo}, the orders of such groups are bounded by some $B$. Since
$G_p$ is residually finite, any $B+1$ distinct elements remain distinct in a single finite quotient, the intersection of
finitely many normal subgroups of finite index. So $|G_p|\leq B$, contradicting the infinitude of $G_p$.
\end{proof}


%======================================================================
\section{Discussion and questions}\label{sec:discussion}

\paragraph{One presentation for every prime.} The algebra of Theorem~A and the groups of Theorem~B come from one
presentation, fixed before the prime. The only place where the prime enters the construction is detection. There,
central elements are separated in finite quotients using long-word phases modulo integers coprime to $p$. The phases
of the hierarchy live modulo products of the encoding lengths $L_j=2^{h_j+1}+1$, and these products may share factors
with $p$. Reducing the phases modulo the part prime to $p$, which is unbounded (Section~\ref{sec:detection}), removes
this dependence. For example, an odd prime divides some $2^m+1$ only if the multiplicative order of $2$ modulo $p$ is
even, so primes such as $7$, $23$, $31$ and $47$ never divide any encoding length.

\paragraph{Exponents.} The groups $G_p$ have unbounded exponent, since a finitely generated residually finite group of
exponent dividing $p^a$ is finite by Zel'manov's theorems~\cite{ZelmanovOdd,ZelmanovTwo}. Whether an infinite
finitely presented group can have finite exponent remains open, and residual finiteness rules out the present method
for that question.

\paragraph{Characteristic zero.} Theorem~C shows that presentations by unsigned word relations, which carry the
positive-characteristic construction, can never give a nil algebra that is not nilpotent in characteristic zero. One
might hope to replace the symmetric sum of Lemma~\ref{lem:symmetric} by an alternating one, using letters that
anticommute. This does not work as stated. Exchanging two whole chunks of lengths $r$ and $s$ then contributes the sign
$(-1)^{rs}$, so chunks of even length still commute and the orbit sums do not become alternating. A characteristic-zero
analogue would need a genuinely different cancellation identity, valid for chunks of arbitrary length, together with an
invariant that keeps arbitrarily long words nonzero.

\paragraph{Questions.}
\begin{enumerate}[label=(\arabic*),leftmargin=*]
\item Is there a finitely presented nil algebra over $\Q$ that is not nilpotent? By Theorem~C its relations must involve
  signs or coefficients.
\item The rank $12$ is inherited from the stability range of the argument and is not optimized. What is the least $n$
  for which $\St_n(R_p)$ has the properties of Theorem~B?
\end{enumerate}

%======================================================================
\section{How this work was done}\label{sec:methods}

\paragraph{Roles.} This paper is the work of its two authors together with AI models from more than one developer.
\begin{itemize}[leftmargin=*]
\item The authors set the direction throughout: they selected the problem from a review of OpenAI's release, fixed the
  scope (every prime, with characteristic zero treated separately), set the verification protocol, decided which lines to
  pursue and which to drop, and take full responsibility for every claim.
\item The models did most of the exploration, derivation, computation and drafting. In particular they found:
  \begin{itemize}
  \item a cancellation bound in characteristic $p$, later found to be a weaker form of Zalesskii's
    theorem~\cite{Zalesskii1985};
  \item the failure of the unordered-frame argument and its ordered-frame replacement;
  \item the sign-corrected group calculations;
  \item the treatment of phases for odd primes;
  \item Theorem~C.
  \end{itemize}
\end{itemize}

\paragraph{Verification.} The extension was worked out in three independent lanes, each with three rounds, and a separate
coordinating instance re-derived the repairs on which the argument rests.
\begin{itemize}[leftmargin=*]
\item \emph{The lanes.} One lane covered the hierarchy, cancellation and phases, and one covered the group-theoretic
  steps. The third was a refutation-first adversary, run by a model from a different developer, which tried to break
  every step and then audited the proofs line by line.
\item \emph{What verification found.} It established that the direct port of the unordered-frame stability argument
  fails over odd fields (Proposition~\ref{prop:unordered-fails}), which led to the ordered-frame replacement.
\item \emph{Exact computation.} Finite identities were checked exactly, in integer arithmetic modulo $p$:
  \begin{itemize}
  \item the symmetric-sum bound of Lemma~\ref{lem:symmetric}, and Zalesskii's degree $dp$, exhaustively on tuples of
    matrix units for small $p$ and $d$;
  \item the matrices of Proposition~\ref{prop:unordered-fails} and of the face comparison in
    Theorem~\ref{thm:stability};
  \item the coordinate changes of Section~\ref{sec:torsion};
  \item the valuation bound of Lemma~\ref{lem:det-pprime}.
  \end{itemize}
  The script, \texttt{verify/checks.py}, is part of the release.
\item \emph{A blind check.} The four load-bearing steps were then checked independently by a model from another developer
  that had not seen our proofs:
  \begin{itemize}
  \item the characteristic-$p$ cancellation and its propagation;
  \item the ordered-frame stability argument;
  \item the ordinary finite presentation with $p$-power root relations;
  \item the detection of central elements.
  \end{itemize}
  It accepted all four, and confirmed by exact computation that the characteristic-two bound of
  \cite[Lemma~3.2]{OAI247periodic} fails over $\F_3$, as Zalesskii's theorem predicts, so a bound depending on $p$ is needed. It
  independently flagged two characteristic-two steps in the imported inputs:
  \begin{itemize}
  \item the permutation lemma~\cite[Lemma~3.3]{OAI247rf};
  \item the swaps in the finite-rank inputs of~\cite[Section~5]{OAI247periodic}.
  \end{itemize}
  Both are replaced in Sections~\ref{sec:stability}--\ref{sec:detection}. It also observed that reducing the phases to
  moduli prime to $p$ removes the need for prime-dependent parameters, which we adopt. It judged Theorem~C correct and
  elementary.
\end{itemize}
The OpenAI manuscripts were read at a fixed commit of their public repository. OpenAI's Lean development formalizes the
characteristic-two periodic group of~\cite{OAI247periodic} and nothing in the present paper. No argument here has been
formally verified; the proofs are given in full, and the checks support them but do not replace them. The
inputs~\cite{OAI247periodic,OAI247rf} were produced by an OpenAI model, as OpenAI states, and have not been refereed.

\paragraph{Priority search.} On 15 October 2026 we searched arXiv, reading math.GR and math.RA newest first and running
full-text queries on finitely presented periodic groups, torsion groups and $p$-groups, nil algebras and Steinberg groups.
We also searched OpenAI's public repository and the general web. We found no infinite finitely presented periodic group,
and no infinite finitely presented residually finite $p$-group, for any odd prime $p$. The nearest recent work is
Ivanov's extension to every odd prime of another construction from the same release, the torsion-free group whose
group algebra over $\F_2$ has zero divisors~\cite{Ivanov2026}. It also replaces a characteristic-two cancellation by a
signed one, but it concerns group algebras and is independent of the present paper. For Theorem~C we found no explicit
statement in the literature on semigroup algebras, and, as said after its proof, we do not claim it as new.


\begin{thebibliography}{99}

\bibitem{BekkaHarpeValette2008}
B.~Bekka, P.~de la Harpe, and A.~Valette,
\emph{Kazhdan's Property $(T)$},
New Mathematical Monographs, vol.~11,
Cambridge University Press, Cambridge, 2008.

\bibitem{Brown1982}
K.~S.~Brown,
\emph{Cohomology of Groups},
Graduate Texts in Mathematics, vol.~87,
Springer-Verlag, New York, 1982.

\bibitem{Burnside1902}
W.~Burnside,
On an unsettled question in the theory of discontinuous groups,
\emph{Quart. J. Pure Appl. Math.} \textbf{33} (1902), 230--238.

\bibitem{ErshovJaikinZapirain2010}
M.~Ershov and A.~Jaikin-Zapirain,
Property $(T)$ for noncommutative universal lattices,
\emph{Invent. Math.} \textbf{179} (2010), no.~2, 303--347.
\url{https://doi.org/10.1007/s00222-009-0218-2}.

\bibitem{Golod1964}
E.~S.~Golod,
On nil-algebras and finitely approximable $p$-groups,
\emph{Izv. Akad. Nauk SSSR Ser. Mat.} \textbf{28} (1964), no.~2, 273--276 (Russian).

\bibitem{GolodShafarevich1964}
E.~S.~Golod and I.~R.~Shafarevich,
On the class field tower,
\emph{Izv. Akad. Nauk SSSR Ser. Mat.} \textbf{28} (1964), no.~2, 261--272 (Russian).

\bibitem{HaettelOsajda2025}
T.~Haettel and D.~Osajda,
Locally elliptic actions, torsion groups, and nonpositively curved spaces,
arXiv:2110.12431v3 (2025).

\bibitem{Ivanov2026}
S.~V.~Ivanov,
Zero divisors in group algebras of torsion-free groups in odd characteristic,
arXiv:2610.17662 (2026).

\bibitem{IvanovPogodaevKanelBelov2025}
I.~A.~Ivanov-Pogodaev and A.~Ya.~Kanel-Belov,
Construction of a nilsemigroup of paths in a countable family of uniformly elliptic complexes,
\emph{Algebra and Logic} \textbf{63} (2025), no.~6, 410--438.
\url{https://doi.org/10.1007/s10469-025-09803-3}.

\bibitem{JonssonWelker2009}
J.~Jonsson and V.~Welker,
Complexes of injective words and their commutation classes,
\emph{Pacific J. Math.} \textbf{243} (2009), no.~2, 313--329.

\bibitem{vanderKallen1980}
W.~van der Kallen,
Homology stability for linear groups,
\emph{Invent. Math.} \textbf{60} (1980), no.~3, 269--295.
\url{https://doi.org/10.1007/BF01390018}.

\bibitem{KrsticMcCool}
S.~Krsti\'c and J.~McCool,
Presenting $\mathrm{GL}_n(k\langle T\rangle)$,
\emph{J. Pure Appl. Algebra} \textbf{141} (1999), no.~2, 175--183.

\bibitem{LyndonSchupp1977}
R.~C.~Lyndon and P.~E.~Schupp,
\emph{Combinatorial Group Theory},
Ergebnisse der Mathematik und ihrer Grenzgebiete, vol.~89,
Springer-Verlag, Berlin, 1977.

\bibitem{Milnor1971}
J.~Milnor,
\emph{Introduction to Algebraic $K$-Theory},
Annals of Mathematics Studies, vol.~72,
Princeton University Press, Princeton, NJ, 1971.

\bibitem{NovikovAdian1968}
P.~S.~Novikov and S.~I.~Adian,
Infinite periodic groups, I--III,
\emph{Math. USSR-Izv.} \textbf{2} (1968), 209--236, 241--479, 665--685.

\bibitem{Okninski1991}
J.~Okni\'nski,
\emph{Semigroup Algebras},
Monographs and Textbooks in Pure and Applied Mathematics \textbf{138},
Marcel Dekker, New York, 1991.

\bibitem{OlshanskiiSapir2003}
A.~Yu.~Ol'shanskii and M.~V.~Sapir,
Non-amenable finitely presented torsion-by-cyclic groups,
\emph{Publ. Math. Inst. Hautes \'Etudes Sci.} \textbf{96} (2003), 43--169.

\bibitem{OAI247periodic}
OpenAI,
An infinite finitely presented periodic group,
OpenAI Math Release preprint, 23 September 2026.
\url{https://github.com/openai/math/blob/main/preprints/An-infinite-finitely-presented-periodic-group-September-23-2026/paper.pdf}.

\bibitem{OAI247rf}
OpenAI,
An infinite finitely presented residually finite 2-group,
OpenAI Math Release preprint, 5 October 2026.
\url{https://github.com/openai/math/blob/main/preprints/An-infinite-finitely-presented-residually-finite-2-group-October-5-2026/residually-finite-torsion.pdf}.

\bibitem{Robinson1996}
D.~J.~S.~Robinson,
\emph{A Course in the Theory of Groups}, 2nd ed.,
Graduate Texts in Mathematics, vol.~80,
Springer-Verlag, New York, 1996.

\bibitem{Schur1904}
I.~Schur,
\"Uber die Darstellung der endlichen Gruppen durch gebrochen lineare Substitutionen,
\emph{J. Reine Angew. Math.} \textbf{127} (1904), 20--50.

\bibitem{Smoktunowicz2008}
A.~Smoktunowicz,
The Jacobson radical of rings with nilpotent homogeneous elements,
\emph{Bull. Lond. Math. Soc.} \textbf{40} (2008), no.~6, 917--928.

\bibitem{Steinberg1967}
R.~Steinberg,
\emph{Lectures on Chevalley Groups},
notes prepared by J.~Faulkner and R.~Wilson,
Yale University, New Haven, 1967.

\bibitem{SuslinTulenbaev1976}
A.~A.~Suslin and M.~S.~Tulenbaev,
Stabilization theorem for the Milnor $K_2$-functor,
\emph{Zap. Nauchn. Sem. Leningrad. Otdel. Mat. Inst. Steklov.} \textbf{64} (1976), 131--152 (Russian);
English translation in \emph{J. Soviet Math.} \textbf{17} (1981), no.~2, 1804--1819.

\bibitem{Valenti1998}
A.~Valenti,
On permanental identities of symmetric and skew-symmetric matrices in characteristic $p$,
\emph{Canad. Math. Bull.} \textbf{41} (1998), no.~1, 118--124.
\url{https://doi.org/10.4153/CMB-1998-018-9}.

\bibitem{Weibel1987}
C.~A.~Weibel,
Module structures on the $K$-theory of graded rings,
\emph{J. Algebra} \textbf{105} (1987), no.~2, 465--483.
\url{https://doi.org/10.1016/0021-8693(87)90210-9}.

\bibitem{WeibelKbook}
C.~A.~Weibel,
\emph{The $K$-Book: An Introduction to Algebraic $K$-Theory},
Graduate Studies in Mathematics, vol.~145,
American Mathematical Society, Providence, RI, 2013.

\bibitem{Zalesskii1985}
A.~E.~Zalesskii,
A symmetric analogue of the Amitsur--Levitzki theorem,
\emph{Vests\={\i} Akad. Navuk BSSR Ser. F\={\i}z.-Mat. Navuk} (1985), no.~4, 108--110 (Russian).

\bibitem{ZelmanovOdd}
E.~I.~Zel'manov,
Solution of the restricted Burnside problem for groups of odd exponent,
\emph{Math. USSR-Izv.} \textbf{36} (1991), no.~1, 41--60.

\bibitem{ZelmanovTwo}
E.~I.~Zel'manov,
A solution of the restricted Burnside problem for $2$-groups,
\emph{Math. USSR-Sb.} \textbf{72} (1992), no.~2, 543--565.

\end{thebibliography}

\end{document}
